
\subsection{Tổng quan kiến trúc}\label{sec:architecture-overview}
\begin{figure}[H]
\centering
\IfFileExists{33_architecture_overview_react_nitro_spring_postgresql.pdf}{
  \includegraphics[width=\textwidth,height=0.76\textheight,keepaspectratio]{33_architecture_overview_react_nitro_spring_postgresql.pdf}
}{
  \IfFileExists{33_architecture_overview_react_nitro_spring_postgresql.png}{
    \includegraphics[width=\textwidth,height=0.76\textheight,keepaspectratio]{33_architecture_overview_react_nitro_spring_postgresql.png}
  }{
    \fbox{\parbox{0.82\linewidth}{Thiếu file \textsf{33\_architecture\_overview\_react\_nitro\_spring\_postgresql.(pdf/png)}.}}
  }
}
\caption{Sơ đồ tổng quan kiến trúc và các trách nhiệm lưu trữ logic của IRMS}
\label{fig:irms-architecture-overview}
\end{figure}

IRMS được thiết kế theo kiến trúc service-based kết hợp event-driven processing, gồm ba lớp chính: giao diện người dùng, dịch vụ nghiệp vụ và dữ liệu cùng các xử lý nền. Kiến trúc này giúp phân tách rõ các chức năng nghiệp vụ và xử lý các tác vụ phụ qua cơ chế sự kiện mà không ảnh hưởng đến luồng giao dịch chính.

Ở lớp giao diện, ReactJS cung cấp nền tảng thống nhất cho các vai trò như lễ tân, phục vụ, bếp, thu ngân và quản lý. Cách tiếp cận này giúp các bộ phận dùng chung cơ chế xác thực, điều hướng và đồng bộ trạng thái, đảm bảo các thông tin như bàn, đơn gọi món, phiếu bếp và hóa đơn luôn được cập nhật nhất quán trong quá trình vận hành.

Phía sau giao diện là \textsf{Frontend Node Server} của TanStack Start/Nitro, chịu trách nhiệm phục vụ React và chuyển tiếp các yêu cầu \textsf{/api} đến \textsf{ApiGatewayApplication}. Backend sử dụng \textsf{GatewayProxyController} làm cổng vào để định tuyến yêu cầu đến các service Spring Boot theo cấu hình \textsf{docker-compose.yml}. Cách tổ chức này giúp frontend không phụ thuộc vào cấu trúc nội bộ của backend, tăng khả năng thay đổi và mở rộng hệ thống.

Lớp backend được xây dựng bằng Spring Boot và tổ chức thành các service theo miền nghiệp vụ. Mỗi service phụ trách một chức năng riêng như \textsf{Reservation \& Seating Service} quản lý đặt bàn và trạng thái bàn, \textsf{Ordering Service} xử lý đơn gọi món, \textsf{Kitchen Coordination Service} điều phối chế biến món ăn, \textsf{Billing \& Settlement Service} quản lý hóa đơn và thanh toán, \textsf{Inventory Monitoring Service} theo dõi tồn kho, \textsf{Reporting Service} phục vụ báo cáo quản trị và \textsf{IAM \/ Audit Service} đảm nhiệm kiểm soát truy cập cùng ghi nhận kiểm toán. Cách phân tách này giúp hệ thống rõ ràng về trách nhiệm, dễ bảo trì và thuận lợi cho mở rộng sau này.

Việc phân chia service theo miền nghiệp vụ giúp kiến trúc phản ánh đúng quy trình vận hành thực tế của nhà hàng. Nhờ ranh giới trách nhiệm rõ ràng, thay đổi ở một nghiệp vụ sẽ ít ảnh hưởng đến các phần khác của hệ thống, từ đó tăng khả năng bảo trì và mở rộng.

Lớp dữ liệu của IRMS sử dụng một PostgreSQL vật lý duy nhất nhưng được phân tách thành các nhóm dữ liệu logic như dữ liệu giao dịch, outbox/inbox, định danh kiểm toán và read model cho báo cáo. Mô hình đọc phục vụ báo cáo được cập nhật qua các xử lý nền thay vì truy vấn trực tiếp lên dữ liệu giao dịch nóng, giúp giảm tải và tách biệt trách nhiệm lưu trữ trong hệ thống.

Ngoài tuyến xử lý chính, kiến trúc còn sử dụng \textsf{outbox\_events}, \textsf{processed\_events}, \textsf{RabbitMQ} và các bộ xử lý nền để thực hiện các tác vụ phụ như gửi thông báo, phát hành biên lai và cập nhật mô hình đọc. Các sự kiện được xử lý bất đồng bộ nhằm giảm tải cho các thao tác trực tiếp như gọi món và thanh toán, đồng thời hỗ trợ cơ chế chống xử lý lặp cho consumer.

Cuối cùng, các điểm tích hợp được đặt ở rìa kiến trúc thông qua các adapter như \textsf{PaymentGateway}, \textsf{NotificationChannelAdapter} và \textsf{ReceiptDeliveryAdapter}. Cách tổ chức này giúp tách biệt lõi nghiệp vụ khỏi thay đổi của hệ thống thanh toán, thông báo và phát hành biên lai, từ đó tăng khả năng bảo trì và mở rộng.

\subsection{Tổng quan cho các góc nhìn chính}\label{sec:architecture-views}

\begin{figure}[H]
\centering
\IfFileExists{21_logical_view_overview_business_capability_layers.pdf}{
  \includegraphics[width=\textwidth,height=0.82\textheight,keepaspectratio]{21_logical_view_overview_business_capability_layers.pdf}
}{
  \IfFileExists{21_logical_view_overview_business_capability_layers.png}{
    \includegraphics[width=\textwidth,height=0.82\textheight,keepaspectratio]{sodologictongquan.png}
  }{
    \fbox{\parbox{0.82\linewidth}{Thiếu file \textsf{21\_logical\_view\_overview\_business\_capability\_layers.(pdf/png)}.}}
  }
}
\caption{Sơ đồ năng lực nghiệp vụ hỗ trợ phân rã module của IRMS}
\label{fig:irms-logical-overview-supplement}
\end{figure}
Hình~\ref{fig:irms-logical-overview-supplement} trình bày sơ đồ năng lực nghiệp vụ của hệ thống IRMS dưới dạng UML Package, trong đó hệ thống được tổ chức dựa trên các năng lực nghiệp vụ (capabilities) thay vì theo giao diện người dùng hoặc cấu trúc kỹ thuật. Mục tiêu của sơ đồ hỗ trợ này là xác định các miền nghiệp vụ độc lập, đồng thời thể hiện mối quan hệ tương tác giữa chúng trước khi được ánh xạ sang các module và service trong các góc nhìn kiến trúc chính.

Trong kiến trúc IRMS, sơ đồ năng lực này được xây dựng theo hướng service-based architecture kết hợp với cơ chế event-driven interaction ở mức khái niệm. Mỗi package tương ứng với một miền năng lực có ranh giới rõ ràng. Sự tương tác giữa các miền được thể hiện thông qua hai cơ chế chính, bao gồm giao tiếp đồng bộ theo mô hình request/response trong luồng nghiệp vụ chính và giao tiếp bất đồng bộ thông qua domain events nhằm truyền tải thông tin vận hành mà không làm tăng mức độ phụ thuộc giữa các miền.

Từ cách tổ chức này, các năng lực nghiệp vụ trong hệ thống được phân nhóm theo vai trò trong kiến trúc tổng thể.

\textbf{Thứ nhất, nhóm năng lực tham chiếu và điều kiện vận hành} đóng vai trò cung cấp dữ liệu nền và đảm bảo điều kiện cho các luồng nghiệp vụ chính. \textit{Menu Management} chịu trách nhiệm quản lý thông tin món ăn, giá bán và trạng thái hiển thị, trong khi \textit{Inventory Monitoring} theo dõi tình trạng tồn kho và phát hiện các điều kiện thiếu hụt nguyên liệu. Nhóm này không trực tiếp tham gia xử lý giao dịch nhưng ảnh hưởng đến tính hợp lệ của dữ liệu trong toàn bộ hệ thống.

\textbf{Thứ hai, chuỗi năng lực cốt lõi} phản ánh toàn bộ vòng đời phục vụ trong nhà hàng. \textit{Reservation \& Seating} đảm nhiệm việc quản lý đặt bàn và phân bổ chỗ ngồi, tạo ngữ cảnh cho \textit{Ordering} trong quá trình tiếp nhận đơn hàng. Sau đó, \textit{Kitchen Coordination} tiếp nhận thông tin từ đơn hàng để điều phối quá trình chế biến món ăn, trong khi \textit{Billing \& Settlement} thực hiện thanh toán và kết thúc phiên phục vụ. Các năng lực trong chuỗi này có quan hệ phụ thuộc logic rõ ràng nhằm đảm bảo luồng nghiệp vụ được xử lý nhất quán.

\textbf{Thứ ba, nhóm năng lực hỗ trợ và phân tích} chịu trách nhiệm xử lý thông tin vận hành theo hướng bất đồng bộ. \textit{Notification Integration} tiếp nhận các sự kiện phát sinh để thực hiện thông báo, trong khi \textit{Reporting \& Analytics} tổng hợp dữ liệu từ nhiều miền nghiệp vụ nhằm phục vụ thống kê và phân tích. Việc tách nhóm này giúp giảm tải cho luồng xử lý chính và tăng khả năng mở rộng hệ thống.

\textbf{Thứ tư, nhóm năng lực kiểm soát hệ thống} đảm bảo các yêu cầu về bảo mật, phân quyền và truy vết. \textit{Identity \& Access Control} quản lý xác thực và phân quyền người dùng, trong khi \textit{Audit \& Compliance} ghi nhận lịch sử hoạt động để phục vụ kiểm toán và tuân thủ. Đây là các năng lực xuyên suốt toàn hệ thống nhưng không tham gia trực tiếp vào luồng nghiệp vụ chính.

\subsubsection{Tổng quan cho góc nhìn mô-đun}
\noindent
Góc nhìn mô-đun mô tả các đơn vị hiện thực chính của backend IRMS, cách chúng được tổ chức thành package, trách nhiệm của từng module và quan hệ phụ thuộc tĩnh giữa các module. Góc nhìn này cũng ghi nhận các cấu trúc dữ liệu dùng chung như \textsf{outbox\_events}, \textsf{processed\_events}, nhật ký kiểm toán và mô hình đọc báo cáo khi chúng ảnh hưởng đến nhiều module. Các cấu trúc dữ liệu này là nhóm bảng hoặc trách nhiệm lưu trữ logic trong cùng PostgreSQL, không phải các database vật lý riêng. Module View không mô tả tiến trình runtime hay luồng message khi hệ thống chạy; các nội dung đó thuộc Component-and-Connector View và Deployment View.

\begin{figure}[H]
\centering
\IfFileExists{22_module_view_overview_layered_dependencies.pdf}{
  \includegraphics[width=\textwidth,height=0.82\textheight,keepaspectratio]{22_module_view_overview_layered_dependencies.pdf}
}{
  \IfFileExists{22_module_view_overview_layered_dependencies.png}{
    \includegraphics[width=\textwidth,height=0.82\textheight,keepaspectratio]{22_module_view_overview_layered_dependencies.png}
  }{
    \fbox{\parbox{0.82\linewidth}{Thiếu file \textsf{22\_module\_view\_overview\_layered\_dependencies.(pdf/png)}.}}
  }
}
\caption{Module View tổng quan theo lớp phụ thuộc của IRMS}
\label{fig:irms-module-overview-supplement}
\end{figure}


Hình~\ref{fig:irms-module-overview-supplement} thể hiện quy tắc phụ thuộc theo lớp bằng Module View, bám vào package và lớp hiện thực trong backend. 

\paragraph{Mô tả}
\begin{itemize}
  \item \textsf{GatewayProxyController} đi qua \textsf{GatewayRouteLocator} và \textsf{RestClientConfiguration} để định tuyến REST, không gọi trực tiếp vào dữ liệu nghiệp vụ.
  \item Các controller như \textsf{OrdersController}, \textsf{KitchenController} và \textsf{BillsController} phụ thuộc vào service ứng dụng tương ứng, đúng với cấu trúc package trong backend.
  \item Các thành phần hạ tầng dùng chung như \textsf{RabbitMqEventPublisher}, \textsf{RabbitMqOutboxRelay} và \textsf{ManualAckConsumerSupport} được đặt riêng để xử lý sự kiện và bộ xử lý sự kiện nền.
\end{itemize}


\paragraph{Catalog mô-đun và ánh xạ hiện thực.}
Bảng~\ref{tab:irms-module-catalog} thể hiện giới mô-đun theo package trong backend. Các hình Module View ở phần này mô tả tổ chức tĩnh và hướng phụ thuộc giữa package, còn các sơ đồ Component-and-Connector phía sau mới mô tả cộng tác runtime.

\begin{table}[H]
\centering
\footnotesize
\renewcommand{\arraystretch}{1.18}
\caption{Catalog mô-đun backend IRMS}
\label{tab:irms-module-catalog}
\begin{tabularx}{\textwidth}{>{\raggedright\arraybackslash}p{0.17\linewidth} >{\raggedright\arraybackslash}X >{\raggedright\arraybackslash}p{0.32\linewidth}}
\toprule
\textbf{Module} & \textbf{Trách nhiệm} \\
\midrule
\textsf{Gateway} & Điểm vào REST, chuẩn hóa đường \textsf{/api/**} và định tuyến đến service nghiệp vụ theo cấu hình.\\
\textsf{Runtime / Bootstrap} & Chứa các entrypoint Spring Boot và profile khởi động từng runtime trong Docker.  \\
\textsf{Reservation} & Đặt bàn, danh sách chờ, sơ đồ bàn, check-in và phiên phục vụ. \\
\textsf{Ordering} & Thực đơn, tạo/xác nhận order, tùy biến món, combo, snapshot giá và gửi bếp. \\
\textsf{Kitchen} & Điều phối phiếu bếp, trạng thái món, deadline, ưu tiên và workflow chế biến. \\
\textsf{Billing} & Lập hóa đơn, tách hóa đơn, thanh toán, hoàn tiền, biên nhận và settlement. \\
\textsf{Inventory} & Theo dõi tồn kho, tiêu hao nguyên liệu, cảnh báo mức tồn thấp và đề xuất nhập hàng.  \\
\textsf{Notification} & Ghi nhận yêu cầu thông báo và chọn adapter kênh gửi nội bộ. \\
\textsf{Reporting} & Dashboard, báo cáo vận hành và các bảng mô hình đọc/projection.  \\
\textsf{Identity/Audit} & Xác thực, vai trò, quyền, session, policy và nhật ký kiểm toán.  \\
\textsf{Common Platform} & Hạ tầng dùng chung: security/context, lỗi chuẩn, workflow, outbox, RabbitMQ và consumer support. \\
\textsf{PDF / Receipt Adapter} & Render/phát hành biên nhận qua adapter nội bộ, không mô tả máy in vật lý hay nhà cung cấp ngoài. \\
\bottomrule
\end{tabularx}
\end{table}

\begin{table}[H]
\centering
\footnotesize
\renewcommand{\arraystretch}{1.18}
\caption{Interface và phụ thuộc chính của các mô-đun}
\label{tab:irms-module-contracts}
\begin{tabularx}{\textwidth}{>{\raggedright\arraybackslash}p{0.16\linewidth} >{\raggedright\arraybackslash}X >{\raggedright\arraybackslash}X}
\toprule
\textbf{Module} & \textbf{Class/interface tiêu biểu và contract nhìn thấy} & \textbf{Phụ thuộc và dữ liệu dùng chung} \\
\midrule
\textsf{Gateway} & \textsf{GatewayProxyController}; \textsf{GatewayRouteLocator}; contract HTTP proxy cho route \textsf{/api/**}. & Dùng cấu hình URL service và common error/context; không gọi trực tiếp dữ liệu nghiệp vụ. \\
\textsf{Runtime / Bootstrap} & \textsf{ApiGatewayApplication}, \textsf{OrderingServiceApplication}, \textsf{KitchenServiceApplication}, \textsf{BillingServiceApplication}, \textsf{ReservationServiceApplication}, \textsf{InventoryServiceApplication}, \textsf{NotificationServiceApplication}, \textsf{ReportingServiceApplication}, \textsf{IdentityAuditServiceApplication}. & Chọn service boundary bằng profile, port và biến môi trường trong \texttt{docker-compose.yml}. \\
\textsf{Reservation} & \textsf{ReservationService}, \textsf{ReservationCheckInService}, \textsf{ReservationTableAssignmentService}. & Dùng identity/common để kiểm soát thao tác; ghi bảng reservation/table session và event qua outbox nếu phát sinh. \\
\textsf{Ordering} & \textsf{OrdersService}, \textsf{OrderMenuSelectionService}, \textsf{ComboPricingPolicy}, \textsf{OrderStatusPolicy}. & Dùng reservation/table-session context, kitchen integration client và common outbox/messaging. \\
\textsf{Kitchen} & \textsf{KitchenService}, \textsf{KitchenOrderRoutingService}, \textsf{KitchenWorkflowMediator}, \textsf{KitchenTicketTransitionPolicy}, \textsf{KitchenTicketItemRepositoryPort}. & Nhận order/kitchen events; dùng \textsf{ManualAckConsumerSupport} và \textsf{processed\_events} cho bộ xử lý sự kiện nền. \\
\textsf{Billing} & \textsf{BillingService}, \textsf{BillingPaymentService}, \textsf{BillingReceiptService}, \textsf{PaymentGateway}, \textsf{ReceiptDeliveryAdapter}, \textsf{BillSplitStrategy}. & Dùng billable snapshot từ ordering, adapter thanh toán/biên nhận và \textsf{DomainEventPublisher}. \\
\textsf{Inventory} & \textsf{InventoryService}, \textsf{InventoryStockChangedConsumer}, \textsf{LowStockSeverityPolicy}, \textsf{ReorderRecommendationStrategy}. & Tiêu thụ event qua RabbitMQ/manual ack; ghi inventory tables và có thể phát cảnh báo tồn thấp. \\
\textsf{Notification} & \textsf{NotificationService}, \textsf{NotificationRequestedConsumer}, \textsf{NotificationChannelAdapter} và các adapter email/SMS/in-app/print. & Nhận yêu cầu thông báo từ sự kiện; kênh gửi được che sau adapter biên nội bộ. \\
\textsf{Reporting} & \textsf{ReportingQueryService}, \textsf{ReportingProjectionConsumer}, các projection/materializer. & Dùng domain events, \textsf{processed\_events} và các bảng \texttt{reporting\_\char`*\_projection}. \\
\textsf{Identity/Audit} & \textsf{AuthService}, \textsf{AuditService}, \textsf{IdentityRepository}, \textsf{IdentityPolicyRepository}, \textsf{AuditRecordingRequestedConsumer}. & Duy trì user/session/role/policy và \textsf{audit\_logs}; các module bảo vệ thao tác dùng chính sách identity. \\
\textsf{Common Platform} & \textsf{TransactionalOutboxPublisher}, \textsf{RabbitMqOutboxRelay}, \textsf{RabbitMqEventPublisher}, \textsf{ManualAckConsumerSupport}, \textsf{ApiEnvelope}. & Crosscutting package cho \textsf{outbox\_events}, \textsf{processed\_events}, envelope/correlation và lỗi chuẩn. \\
\bottomrule
\end{tabularx}
\end{table}

\begin{table}[H]
\centering
\scriptsize
\renewcommand{\arraystretch}{1.15}
\setlength{\tabcolsep}{2pt}
\caption{Các loại quan hệ dùng trong Module View}
\label{tab:irms-module-relation-types}
\begin{tabularx}{\textwidth}{>{\raggedright\arraybackslash}p{0.18\linewidth} >{\raggedright\arraybackslash}X >{\raggedright\arraybackslash}p{0.34\linewidth}}
\toprule
\textbf{Relation type} & \textbf{Ý nghĩa trong Module View} & \textbf{Ví dụ trong IRMS} \\
\midrule
\textsf{is-part-of} & Module/package con thuộc một package lớn hơn. & \textsf{ordering.api}, \textsf{ordering.application}, \textsf{ordering.domain}, \textsf{ordering.infrastructure} thuộc \textsf{SA.irms.ordering}. \\
\textsf{uses / depends-on} & Phụ thuộc source-level qua import, constructor, method call, config hoặc bean wiring. & \textsf{OrdersController} phụ thuộc \textsf{OrdersService}; \textsf{BillingPaymentService} phụ thuộc \textsf{PaymentGateway}. \\
\textsf{allowed-to-use} & Quy tắc tầng được phép dùng tầng thấp hơn hoặc package dùng chung. & \textsf{*.api} dùng \textsf{*.application}; \textsf{*.application} dùng \textsf{*.domain} và port/interface. \\
\textsf{data dependency} & Module đọc/ghi hoặc bị ảnh hưởng bởi cấu trúc dữ liệu dùng chung. & \textsf{RabbitMqOutboxRelay} dùng \textsf{outbox\_events}; bộ xử lý sự kiện nền dùng \textsf{processed\_events}. \\
\textsf{crosscuts} & Package dùng chung hỗ trợ nhiều module nghiệp vụ. & \textsf{SA.irms.common} cung cấp security/context, outbox, messaging, workflow và lỗi chuẩn. \\
\textsf{is-a} & Quan hệ kế thừa hoặc hiện thực interface ở mức class/interface. & \textsf{CashPaymentGateway} hiện thực \textsf{PaymentGateway}; notification adapters hiện thực \textsf{NotificationChannelAdapter}. \\
\bottomrule
\end{tabularx}
\end{table}

\begin{table}[H]
\centering
\footnotesize
\renewcommand{\arraystretch}{1.15}
\caption{Cấu trúc dữ liệu dùng chung trong Module View}
\label{tab:irms-global-data-structures}
\begingroup
\small
\setlength{\tabcolsep}{3pt}
\renewcommand{\arraystretch}{1.18}

\begin{tabularx}{\textwidth}{@{}
  >{\raggedright\arraybackslash}p{0.18\textwidth}
  >{\raggedright\arraybackslash}p{0.22\textwidth}
  >{\raggedright\arraybackslash}X
  >{\raggedright\arraybackslash}p{0.23\textwidth}
@{}}
\toprule
\textbf{Cấu trúc} &
\textbf{Module duy trì} &
\textbf{Vai trò trong góc nhìn mô-đun} &
\textbf{Thành phần liên quan} \\
\midrule

\textsf{outbox\_events} &
\textsf{SA.irms.common.outbox} và các module phát event. &
Bảng dùng chung để ghi sự kiện cần phát bất đồng bộ trước khi chuyển sang \textsf{RabbitMQ}. &
\begin{tabular}[t]{@{}l@{}}
\textsf{TransactionalOutbox}\\
\textsf{Publisher};\\
\textsf{RabbitMqOutboxRelay}.
\end{tabular} \\

\textsf{processed\_events} &
\textsf{SA.irms.common.messaging}. &
Bảng ghi nhận sự kiện đã xử lý, giúp tránh xử lý lặp trong các luồng nền. &
\begin{tabular}[t]{@{}l@{}}
\textsf{ManualAckConsumer}\\
\textsf{Support}.
\end{tabular} \\

\begin{tabular}[t]{@{}l@{}}
\texttt{reporting\_\char`*\_}\\
\texttt{projection}
\end{tabular} &
\textsf{SA.irms.reporting}. &
Mô hình đọc cho dashboard và báo cáo, được cập nhật qua các thành phần tổng hợp dữ liệu. &
\begin{tabular}[t]{@{}l@{}}
\textsf{ReportingProjection}\\
\textsf{Consumer};\\
\textsf{ReportingQueryService}.
\end{tabular} \\

\begin{tabular}[t]{@{}l@{}}
\textsf{audit\_logs},\\
\textsf{audit\_details}
\end{tabular} &
\textsf{SA.irms.identity}. &
Dữ liệu kiểm toán cho thao tác nhạy cảm và truy vết vận hành. &
\begin{tabular}[t]{@{}l@{}}
\textsf{AuditService};\\
\textsf{AuditRecording}\\
\textsf{RequestedConsumer}.
\end{tabular} \\

\begin{tabular}[t]{@{}l@{}}
\textsf{identity/session/}\\
\textsf{policy data}
\end{tabular} &
\textsf{SA.irms.identity}. &
Dữ liệu xác thực, phân quyền, session và chính sách được nhiều module tham chiếu. &
\begin{tabular}[t]{@{}l@{}}
\textsf{AuthService};\\
\textsf{IdentityPolicy}\\
\textsf{Repository}.
\end{tabular} \\

\bottomrule
\end{tabularx}
\endgroup
\end{table}

\subsubsection{Tổng quan cho góc nhìn thành phần và kết nối}

Góc nhìn thành phần và kết nối mô tả các thành phần có mặt ở runtime và cách chúng giao tiếp qua REST, JDBC, AMQP, outbox và các adapter biên. Khác với Module View, góc nhìn này không chỉ mô tả package tĩnh mà làm rõ component nào cung cấp interface, component nào yêu cầu interface, port nào nằm ở biên subsystem và connector nào lắp ghép các interface đó. Trong IRMS, trọng tâm của C\&C View là tuyến đi từ frontend/API Gateway vào controller/service nghiệp vụ, tuyến ghi dữ liệu vào PostgreSQL vật lý duy nhất, và tuyến event-driven đi qua outbox, RabbitMQ, bộ xử lý sự kiện nền, read model, notification và audit.

Phần chi tiết của Component-and-Connector View được trình bày ở Mục~\ref{sec:main-components} và các hình thành phần bên dưới. Các hình này cần được đọc như sơ đồ runtime collaboration có component, port/interface và connector, không phải package diagram hay class diagram.

\subsubsection{Tổng quan cho góc nhìn phân bổ}
\begin{figure}[H]
\centering
\IfFileExists{23_allocation_view_overview_software_to_runtime_mapping.pdf}{
  \includegraphics[width=\textwidth,height=0.84\textheight,keepaspectratio]{23_allocation_view_overview_software_to_runtime_mapping.pdf}
}{
  \IfFileExists{23_allocation_view_overview_software_to_runtime_mapping.png}{
    \includegraphics[width=\textwidth,height=0.84\textheight,keepaspectratio]{23_allocation_view_overview_software_to_runtime_mapping.png}
  }{
    \fbox{\parbox{0.82\linewidth}{Thiếu file \textsf{23\_allocation\_view\_overview\_software\_to\_runtime\_mapping.(pdf/png)}.}}
  }
}
\caption{Sơ đồ tổng quan cho góc nhìn phân bổ của IRMS}
\label{fig:irms-allocation-overview}
\end{figure}

\paragraph{Mô tả:}
\begin{itemize}
  \item \textbf{Lớp biên và thiết bị}: là nơi phát sinh thao tác nghiệp vụ thực tế như máy của lễ tân, máy của phục vụ, màn hình bếp, máy thu ngân và màn hình quản lý.
  \item \textbf{Lớp thực thi}: gồm \textsf{frontend container}, \textsf{api-gateway container}, nhóm container service Spring Boot và \textsf{irms-migration container}; đây là nơi các năng lực được hiện thực thành tiến trình hoặc vùng thực thi cụ thể.
  \item \textbf{Lớp dữ liệu và sự kiện}: gồm \textsf{PostgreSQL container}, \textsf{RabbitMQ container}, \textsf{outbox\_events}, \textsf{processed\_events} và các bảng hoặc view báo cáo; chúng tách dữ liệu giao dịch, dữ liệu đọc tổng hợp và hàng đợi sự kiện.
  \item \textbf{Lớp hệ thống ngoài}: hình này không vẽ hệ thống ngoài độc lập cho thanh toán, thông báo hoặc máy in vật lý; các năng lực đó được đặt sau adapter biên nội bộ trong phạm vi hiện tại.
\end{itemize}


\subsection{Đặc trưng kiến trúc}
Những \emph{đặc trưng kiến trúc} quan trọng là:
\begin{enumerate}
  \item \textbf{Khả năng đáp ứng}: màn hình gọi món, màn hình bếp và màn hình thu ngân phải phản hồi nhanh để không làm chậm tuyến phục vụ trực tiếp với khách.
  \item \textbf{Độ tin cậy}: đơn gọi món, thanh toán, hoàn tiền và cập nhật trạng thái bàn là các luồng không được mất dữ liệu hoặc xử lý lặp sai lệch.
  \item \textbf{Khả năng mở rộng}: điểm nóng thường dồn vào giờ cao điểm và tập trung ở các miền \textsf{Ordering}, \textsf{Kitchen Coordination} và \textsf{Billing} thay vì toàn hệ thống.
  \item \textbf{Bảo mật và khả năng kiểm toán}: mọi thao tác nhạy cảm phải được kiểm soát bằng quyền và truy vết lại được về sau.
  \item \textbf{Khả năng quan sát}: hệ thống phải đủ nhật ký, vết thực thi và số đo để thấy độ trễ của phiếu bếp, thanh toán và các điểm nghẽn vận hành.
  \item \textbf{Khả năng sửa đổi}: thực đơn, giá, khuyến mãi, ngưỡng cảnh báo tồn kho, kênh gửi thông báo và bộ kết nối thanh toán thay đổi thường xuyên nên cần biên thay đổi hẹp.
\end{enumerate}

\subsection{Sơ đồ năng lực nghiệp vụ hỗ trợ}

\subsubsection{Sơ đồ phân rã năng lực logic}

\begin{figure}[H]
\centering
\IfFileExists{04_initial_logical_components_cua_irms.pdf}{
  \includegraphics[width=\textwidth,height=0.78\textheight,keepaspectratio]{04_initial_logical_components_cua_irms.pdf}
}{
  \IfFileExists{04_initial_logical_components_cua_irms.png}{
    \includegraphics[width=\textwidth,height=0.78\textheight,keepaspectratio]{04_initial_logical_components_cua_irms.png}
  }{
    \fbox{\parbox{0.82\linewidth}{Thiếu file \texttt{04\_initial\_logical\_components\_cua\_irms.(pdf/png)}.}}
  }
}
\caption{Sơ đồ phân rã năng lực logic}
\label{fig:irms-initial-components}
\end{figure}

Hình~\ref{fig:irms-initial-components} thể hiện bước phân rã ban đầu các năng lực nghiệp vụ trong hệ thống IRMS.

Trong sơ đồ này, các năng lực được xác định theo hướng chức năng bề mặt như quản lý bàn, tiếp nhận đơn gọi món, điều phối bếp, thanh toán, quản lý tồn kho, báo cáo, kiểm soát truy cập và giám sát hệ thống. 

Mục tiêu là nhận diện đầy đủ phạm vi năng lực của hệ thống IRMS. Các năng lực được mô tả theo cách gần với quy trình vận hành thực tế để dễ kiểm tra và tránh bỏ sót chức năng quan trọng.

Sơ đồ này được xem là bước khởi tạo trong quá trình phân tích logic, đóng vai trò nền tảng để tiếp tục tái cấu trúc và gom nhóm các năng lực thành các miền nghiệp vụ ổn định hơn ở bước tiếp theo.
\subsubsection{Sơ đồ phân rã năng lực logic hoàn chỉnh}

\begin{figure}[H]
\centering
\IfFileExists{05_refined_logical_components_cua_irms.pdf}{
  \includegraphics[width=\textwidth,height=0.78\textheight,keepaspectratio]{05_refined_logical_components_cua_irms.pdf}
}{
  \IfFileExists{05_refined_logical_components_cua_irms.png}{
    \includegraphics[width=\textwidth,height=0.78\textheight,keepaspectratio]{05_refined_logical_components_cua_irms.png}
  }{
    \fbox{\parbox{0.82\linewidth}{Thiếu file \texttt{05\_refined\_logical\_components\_cua\_irms.(pdf/png)}.}}
  }
}
\caption{Sơ đồ phân rã năng lực logic của IRMS}
\label{fig:irms-refined-components}
\end{figure}
Trong cấu trúc này, các chức năng được tổ chức theo các miền nghiệp vụ có tính kết dính cao và phản ánh đầy đủ quy trình vận hành của nhà hàng, bao gồm:
\begin{itemize}
    \item Table \& Reservation
    \item Menu Management
    \item Ordering
    \item Kitchen Coordination
    \item Billing \& Settlement
    \item Inventory Monitoring
    \item Reporting \& Analytics
    \item Identity \& Access Control
    \item Audit \& Compliance
    \item Notification Integration
\end{itemize}

Việc tái cấu trúc này giúp chuyển từ cách nhìn theo chức năng thao tác người dùng sang cách nhìn theo miền trách nhiệm nghiệp vụ. Thay vì xem hệ thống như tập hợp các chức năng rời rạc, mỗi năng lực được định nghĩa như một miền logic độc lập, có ranh giới rõ ràng và chịu trách nhiệm cho một phần xác định trong vòng đời vận hành của nhà hàng.

Từ góc độ kiến trúc, đây là bước chuyển từ functional decomposition sang domain-oriented decomposition. Các năng lực lúc này không chỉ mang ý nghĩa chức năng, mà trở thành các đơn vị nghiệp vụ có tính độc lập tương đối, có thể được ánh xạ trực tiếp hoặc gián tiếp sang các service trong kiến trúc service-based.

Bên cạnh đó, cấu trúc này cũng tạo nền tảng cho mô hình event-driven processing trong các góc nhìn kiến trúc tiếp theo. Các miền năng lực có thể trao đổi thông tin thông qua sự kiện nghiệp vụ thay vì phụ thuộc trực tiếp lẫn nhau, từ đó giảm coupling giữa các thành phần và tăng khả năng mở rộng của hệ thống trong dài hạn.
\subsubsection{Sơ đồ chuỗi phục vụ}

\begin{figure}[H]
\centering
\IfFileExists{34_logical_view_core_service_chain.pdf}{
  \includegraphics[width=\textwidth,height=0.84\textheight,keepaspectratio]{34_logical_view_core_service_chain.pdf}
}{
  \IfFileExists{34_logical_view_core_service_chain.png}{
    \includegraphics[width=\textwidth,height=0.84\textheight,keepaspectratio]{34_logical_view_core_service_chain.png}
  }{
    \fbox{\parbox{0.82\linewidth}{Thiếu file \texttt{34\_logical\_view\_core\_service\_chain.(pdf/png)}.}}
  }
}
\caption{Sơ đồ năng lực cho chuỗi phục vụ lõi của IRMS}
\label{fig:irms-logical-core-chain}
\end{figure}
Hình~\ref{fig:irms-logical-core-chain} mô tả cách các năng lực logic lõi tương tác với nhau trong chuỗi phục vụ chính của hệ thống.

Chuỗi này bao gồm bốn miền năng lực chính:

\begin{itemize}
    \item Reservation \& Seating: thiết lập và duy trì ngữ cảnh phục vụ
    \item Ordering: ghi nhận và quản lý đơn gọi món
    \item Kitchen Coordination: điều phối quá trình chế biến món ăn
    \item Billing \& Settlement: xử lý thanh toán và kết thúc phiên phục vụ
\end{itemize}

Sơ đồ thể hiện rằng hệ thống không vận hành như tập hợp các chức năng độc lập, mà được tổ chức theo một chuỗi năng lực có thứ tự rõ ràng, phản ánh trực tiếp vòng đời phục vụ của nhà hàng. Mỗi năng lực tương ứng với một giai đoạn trong quy trình nghiệp vụ tổng thể và chỉ chịu trách nhiệm trong phạm vi của giai đoạn đó.

Trong chuỗi này, các miền nghiệp vụ chỉ trao đổi lượng thông tin tối thiểu cần thiết để đảm bảo tính nhất quán của luồng xử lý. Reservation \& Seating cung cấp ngữ cảnh phiên bàn hợp lệ; Ordering tạo và xác nhận thông tin đơn gọi món; Kitchen Coordination tiếp nhận và xử lý trạng thái chế biến; trong khi Billing \& Settlement sử dụng dữ liệu đã được chốt để thực hiện quyết toán và kết thúc phiên phục vụ.

Cách tổ chức này giúp đảm bảo luồng nghiệp vụ chính có tính tuyến tính và dễ kiểm soát, đồng thời giảm thiểu sự phụ thuộc trực tiếp giữa các miền nghiệp vụ. Về mặt kiến trúc, đây là cơ sở để hiện thực hóa các tương tác giữa các năng lực thông qua cơ chế gọi dịch vụ đồng bộ trong luồng chính và cơ chế sự kiện nghiệp vụ (domain events) cho các tương tác mở rộng trong các góc nhìn triển khai phía sau.

\subsection{Góc nhìn phân rã}\label{sec:module-view}
Ở góc nhìn phân rã, hệ thống được chia thành hai lớp lớn:
\textbf{các ứng dụng phía người dùng} và \textbf{các service nghiệp vụ }.

\subsubsection{Các ứng dụng phía người dùng}
\begin{itemize}
  \item \textbf{Host/Reservation UI}: phục vụ đặt bàn, danh sách chờ, sơ đồ bàn và check-in.
  \item \textbf{Server POS UI}: phục vụ gọi món tại bàn, thay đổi món, theo dõi phiếu bếp và thực hiện tính tiền.
  \item \textbf{Kitchen Display UI}: hiển thị phiếu bếp theo trạm, mức ưu tiên, thời gian dự kiến hoàn thành và tiến độ chế biến.
  \item \textbf{Cashier UI}: tạo hóa đơn, tách hóa đơn, thanh toán, biên lai và hoàn tiền.
  \item \textbf{Manager Portal}: quản lý thực đơn, chương trình khuyến mãi, cảnh báo tồn kho, lịch làm việc và báo cáo.
\end{itemize}

\subsubsection{Các dịch vụ nghiệp vụ }
\begin{itemize}
  \item \textbf{Reservation \& Seating Service}: quản lý đặt bàn, danh sách chờ, bàn ăn và phiên phục vụ tại bàn.
  \item \textbf{Ordering Service}: quản lý thực đơn, tùy chọn món, đơn gọi món, dòng món và trạng thái của đơn từ đầu đến cuối.
  \item \textbf{Kitchen Coordination Service}: điều phối phiếu bếp, trạm bếp và cam kết thời gian ra món.
  \item \textbf{Billing Service}: quản lý hóa đơn, tách hóa đơn, thanh toán, biên lai, hoàn tiền và áp dụng khuyến mãi.
  \item \textbf{Inventory Service}: quản lý công thức, mặt hàng tồn kho, giao dịch kho và quy tắc cảnh báo mức tồn thấp.
  \item \textbf{Reporting Service}: quản lý mô hình đọc, ảnh chụp dữ liệu báo cáo và xuất báo cáo.
  \item \textbf{IAM/Access Service}: định danh, vai trò, quyền và thực thi chính sách truy cập.
\end{itemize}


\subsubsection{Sơ đồ tổng quan}
\begin{figure}[H]
\centering
\IfFileExists{06_module_view_overall_irms.pdf}{
  \includegraphics[width=\textwidth,height=0.82\textheight,keepaspectratio]{06_module_view_overall_irms.pdf}
}{
  \IfFileExists{06_module_view_overall_irms.png}{
    \includegraphics[width=\textwidth,height=0.82\textheight,keepaspectratio]{06_module_view_overall_irms.png}
  }{
    \fbox{\parbox{0.82\linewidth}{Thiếu file \textsf{06\_module\_view\_overall\_irms.(pdf/png)}.}}
  }
}
\caption{Module View tổng quan theo package backend của IRMS}
\label{fig:irms-module-overall}
\end{figure}

Sơ đồ thể hiện \textbf{cách các package và class trung tâm được chia theo nghĩa nghiệp vụ}. \textsf{ReservationService} đại diện cho vùng đặt bàn; \textsf{OrdersService} xử lý tạo và xác nhận đơn; \textsf{KitchenService} xử lý trạng thái phiếu bếp; \textsf{BillingService} xử lý lập hóa đơn; còn \textsf{InventoryService}, \textsf{ReportingQueryService}, \textsf{NotificationService}, \textsf{AuthService} và \textsf{AuditService} giữ các vùng hỗ trợ và kiểm soát.

Một điểm quan trọng khác của sơ đồ là \textbf{các phụ thuộc giữa thành phần}. \textsf{apiFetch} gọi các endpoint đã khai báo trong \textsf{ENDPOINTS}, sau đó \textsf{GatewayProxyController} định tuyến đến service phù hợp. Các service nghiệp vụ cùng phát sự kiện qua \textsf{RabbitMqEventPublisher}, còn \textsf{RabbitMqOutboxRelay} xử lý phần chuyển tiếp nền. 

\paragraph{Mô tả}
\begin{itemize}
  \item \textbf{Các thành phần frontend và gateway}: \textsf{apiFetch}, \textsf{ENDPOINTS}, \textsf{permissions.ts}, \textsf{GatewayProxyController} và \textsf{GatewayRouteLocator}.
  \item \textbf{Các service ứng dụng}: \textsf{ReservationService}, \textsf{OrdersService}, \textsf{KitchenService}, \textsf{BillingService}, \textsf{InventoryService}, \textsf{ReportingQueryService}, \textsf{NotificationService}, \textsf{AuthService} và \textsf{AuditService}.
  \item \textbf{Các thành phần sự kiện dùng chung}: \textsf{RabbitMqEventPublisher} và \textsf{RabbitMqOutboxRelay}.
\end{itemize}



\begin{figure}[H]
\centering
\IfFileExists{07_module_view_core_fulfillment_path.pdf}{
  \includegraphics[width=0.97\linewidth,height=0.84\textheight,keepaspectratio]{07_module_view_core_fulfillment_path.pdf}
}{
  \IfFileExists{07_module_view_core_fulfillment_path.png}{
    \includegraphics[width=0.97\linewidth,height=0.84\textheight,keepaspectratio]{07_module_view_core_fulfillment_path.png}
  }{
    \fbox{\parbox{0.82\linewidth}{Thiếu file \textsf{07\_module\_view\_core\_fulfillment\_path.(pdf/png)}.}}
  }
}
\caption{Module View chi tiết cho tuyến xử lý nghiệp vụ lõi}
\label{fig:irms-module-core}
\end{figure}



Trong cấu trúc này, \textsf{OrdersService} được đặt ở trung tâm để phản ánh vai trò cốt lõi của nghiệp vụ gọi món, nơi hệ thống phải kiểm tra trạng thái bàn, tính hợp lệ của món và tùy chọn, đồng thời chuyển thông tin sang bếp thông qua \textsf{RemoteKitchenOrderRoutingClient}. Service này tập trung vào xử lý dữ liệu gọi món và điều phối nghiệp vụ trực tiếp, tránh trở thành nơi chứa toàn bộ logic như khuyến mãi, tồn kho, thông báo hay phân tích.

Mối quan hệ giữa \textsf{OrdersService}, \textsf{KitchenOrderRoutingService} và \textsf{BillingBillCreationService} cũng thể hiện luồng nghiệp vụ rõ ràng: từ xác nhận đơn gọi món, sang điều phối chế biến tại bếp, rồi đến lập hóa đơn và thanh toán. Các service phía sau sử dụng dữ liệu đáng tin cậy từ đơn gọi món nhưng không can thiệp ngược vào cấu trúc nghiệp vụ gốc.
\paragraph{Mô tả}
\begin{itemize}
  \item \textbf{\textsf{ReservationCheckInService} và \textsf{ReservationTableAssignmentService}}: cung cấp ngữ cảnh bàn đang phục vụ.
  \item \textbf{\textsf{OrdersService}}: giữ vai trò trục trung tâm của tuyến nghiệp vụ nóng.
  \item \textbf{\textsf{KitchenOrderRoutingService} và \textsf{BillingBillCreationService}}: là hai hướng lan truyền quan trọng sau khi đơn được xác nhận.
\end{itemize}



\begin{figure}[H]
\centering
\IfFileExists{08_module_view_support_and_cross_cutting_modules.pdf}{
  \includegraphics[width=0.97\linewidth,height=0.84\textheight,keepaspectratio]{08_module_view_support_and_cross_cutting_modules.pdf}
}{
  \IfFileExists{08_module_view_support_and_cross_cutting_modules.png}{
    \includegraphics[width=0.97\linewidth,height=0.84\textheight,keepaspectratio]{08_module_view_support_and_cross_cutting_modules.png}
  }{
    \fbox{\parbox{0.82\linewidth}{Thiếu file \textsf{08\_module\_view\_support\_and\_cross\_cutting\_modules.(pdf/png)}.}}
  }
}
\caption{Module View cho các package hỗ trợ và xuyên suốt}
\label{fig:irms-module-support}
\end{figure}

Hình trên mô tả các package hỗ trợ và các thành phần xử lý nền của hệ thống. Đây là những thành phần không nằm trực tiếp trong các giao dịch chính như gọi món hay thanh toán, nhưng vẫn đóng vai trò quan trọng đối với vận hành và quản lý hệ thống. Góc nhìn này tập trung vào các module liên quan đến messaging, event processing và runtime support thay vì các lớp nghiệp vụ lõi.

Trong kiến trúc IRMS, các consumer như \textsf{InventoryStockChangedConsumer}, \textsf{ReportingProjectionConsumer}, \textsf{NotificationRequestedConsumer} và \textsf{AuditRecordingRequestedConsumer} đảm nhận các chức năng khác nhau. \textsf{InventoryStockChangedConsumer} cập nhật biến động tồn kho dựa trên dữ liệu vận hành; \textsf{ReportingProjectionConsumer} xây dựng dữ liệu phục vụ báo cáo; \textsf{NotificationRequestedConsumer} xử lý việc gửi thông báo; còn \textsf{AuditRecordingRequestedConsumer} ghi nhận dữ liệu kiểm toán và theo dõi thao tác trong hệ thống.

Các thành phần này nhận dữ liệu sự kiện từ các service nghiệp vụ như \textsf{ReservationService}, \textsf{OrdersService}, \textsf{KitchenService}, \textsf{BillingService} và \textsf{InventoryService} thông qua \textsf{RabbitMqEventPublisher}. Nhờ đó, các chức năng hỗ trợ có thể được xử lý bất đồng bộ mà không ảnh hưởng trực tiếp đến các giao dịch chính.

Việc tách riêng các service hỗ trợ giúp hệ thống dễ mở rộng và thuận tiện hơn cho quản trị vận hành. Các chức năng như báo cáo, cảnh báo tồn kho, phân quyền hay kiểm toán có thể hoạt động độc lập với tuyến xử lý chính của hệ thống.

Tuy nhiên, cách tổ chức này cũng dẫn đến độ trễ nhất định trong các chức năng hỗ trợ. Ví dụ, dữ liệu báo cáo hoặc cảnh báo tồn kho có thể cập nhật chậm hơn vài giây so với giao dịch thực tế. Đánh đổi này được chấp nhận vì các thao tác cần phản hồi ngay trong nhà hàng vẫn là đặt bàn, gọi món, chế biến và thanh toán.
\paragraph{Các thành phần chính trong hình.}
\begin{itemize}
  \item \textbf{\textsf{ReservationService}, \textsf{OrdersService}, \textsf{KitchenService}, \textsf{BillingService} và \textsf{InventoryService}}: là các thành phần nguồn phát sinh dữ kiện vận hành để nhóm hỗ trợ tiêu thụ.
  \item \textbf{\textsf{RabbitMqEventPublisher} và \textsf{ManualAckConsumerSupport}}: là lớp trung gian chuẩn hóa và xử lý sự kiện trước khi sang nhóm hỗ trợ.
  \item \textbf{\textsf{InventoryStockChangedConsumer}, \textsf{ReportingProjectionConsumer} và \textsf{NotificationRequestedConsumer}}: đại diện cho nhóm suy diễn, tổng hợp và tích hợp có thể chạy lệch nhịp với đường nóng.
  \item \textbf{\textsf{AuditRecordingRequestedConsumer} và \textsf{AuditService}}: đại diện cho nhóm kiểm soát xuyên suốt.
  \item \textbf{Các adapter kênh thông báo}: gồm \textsf{EmailNotificationChannelAdapter}, \textsf{SmsNotificationChannelAdapter}, \textsf{InAppNotificationChannelAdapter} và \textsf{PrintNotificationChannelAdapter}.
\end{itemize}


\subsubsection{Chi tiết Module View cho tuyến order-to-cash}
\begin{figure}[H]
\centering
\IfFileExists{35_module_view_order_to_cash_slices.pdf}{
  \includegraphics[width=\textwidth,height=0.84\textheight,keepaspectratio]{35_module_view_order_to_cash_slices.pdf}
}{
  \IfFileExists{35_module_view_order_to_cash_slices.png}{
    \includegraphics[width=\textwidth,height=0.84\textheight,keepaspectratio]{35_module_view_order_to_cash_slices.png}
  }{
    \fbox{\parbox{0.82\linewidth}{Thiếu file \textsf{35\_module\_view\_order\_to\_cash\_slices.(pdf/png)}.}}
  }
}
\caption{Module View chi tiết cho tuyến order-to-cash của IRMS}
\label{fig:irms-module-order-to-cash}
\end{figure}

Hình~\ref{fig:irms-module-order-to-cash} đi sâu vào luồng nghiệp vụ Order-to-Cash: từ lúc khách nhận bàn, gọi món, báo xuống bếp, thanh toán cho đến khi hệ thống phát sự kiện thông báo và cập nhật báo cáo. Dưới góc nhìn Module View cho riêng lát cắt nghiệp vụ này, sơ đồ tập trung làm rõ các package và class chính tham gia xử lý như \textsf{ReservationCheckInService}, \textsf{OrdersService}, \textsf{KitchenOrderRoutingService}, \textsf{BillingPaymentService} và \textsf{RabbitMqEventPublisher}.

Trong thiết kế này, \textbf{mỗi service sẽ tự quản lý lớp điều phối và các port hạ tầng của riêng nó}. Cụ thể, \textsf{ReservationSeatingMediator} lo việc nhận và xếp bàn; \textsf{OrdersService} xử lý tạo đơn; \textsf{RemoteKitchenOrderRoutingClient} sẽ giao tiếp với \textsf{KitchenIntegrationController}; tiến độ làm món được \textsf{KitchenTicketItemWorkflowService} theo dõi; còn \textsf{BillingPaymentService} thì kết nối qua interface \textsf{PaymentGateway}. Các service này hoạt động độc lập, chỉ phụ thuộc nhau qua API, event hoặc interface chứ không truy cập chéo vào cấu trúc nội bộ của nhau.

Sơ đồ cũng thể hiện rõ việc phân tách giữa \textbf{luồng xử lý đồng bộ} (synchronous) và \textbf{luồng xử lý bất đồng bộ qua event} (asynchronous). Các thao tác cần phản hồi cho user ngay lập tức (nhận bàn, chốt order, lập hóa đơn, thanh toán) sẽ được chạy đồng bộ. Ngược lại, các tác vụ "hậu cần" sẽ được đẩy vào \textsf{RabbitMqEventPublisher} để các consumer (\textsf{ReportingProjectionConsumer}, \textsf{NotificationRequestedConsumer}, \textsf{AuditRecordingRequestedConsumer}) tự động bắt lấy và xử lý. Việc tách bạch này giúp luồng chính chạy nhanh hơn và không bị block bởi các tác vụ phụ như gửi thông báo, lưu log audit hay tổng hợp báo cáo.

Cuối cùng, các controller và gateway đóng vai trò là các điểm chạm (entry point) đứng trước các service nghiệp vụ. Đây là nơi tiếp nhận request, validate dữ liệu và chuẩn hóa lời gọi. Nhờ có tầng biên này, các service nghiệp vụ bên trong không phải gánh thêm phần logic kỹ thuật, đồng thời đảm bảo ranh giới giao dịch (transaction boundary) của từng service được giữ nguyên vẹn.
\paragraph{Mô tả}
\begin{itemize}
  \item \textbf{Lát cắt service} được thể hiện bằng package và class thật: controller, service ứng dụng, interface và client tích hợp.
  \item \textbf{Tuyến phụ thuộc tĩnh} chỉ đi qua giao diện công khai và client REST cần thiết, nhờ đó tránh phụ thuộc chéo vào chi tiết nội bộ.
  \item \textbf{Tuyến hệ quả sự kiện} được tách sang \textsf{RabbitMqEventPublisher} và các bộ xử lý sự kiện nền, bảo vệ đường nóng khỏi tác vụ phụ.
  \item \textbf{PaymentGateway} là interface với hai hiện thực \textsf{CashPaymentGateway} và \textsf{ExternalReferencePaymentGateway}.
  \item \textbf{Quy tắc phụ thuộc runtime} ở hình này là cơ sở để kiểm chứng phần hiện thực hiện thực: nếu service truy cập xuyên qua chi tiết nội bộ của nhau, kiến trúc sẽ bị phá vỡ ngay ở mức thiết kế.
\end{itemize}


\subsection{Góc nhìn thành phần và kết nối}\label{sec:component-connector-view}

Góc nhìn này cụ thể hóa các quyết định về gateway boundary, REST đồng bộ, event-driven processing, outbox, RabbitMQ và idempotency trong \adrref{ADR-0003}{adr:0003-gateway-boundary}, \adrref{ADR-0004}{adr:0004-rest-sync}, \adrref{ADR-0005}{adr:0005-event-driven}, \adrref{ADR-0006}{adr:0006-outbox}, \adrref{ADR-0007}{adr:0007-rabbitmq} và \adrref{ADR-0010}{adr:0010-idempotency}.
Ở góc nhìn thành phần và kết nối, hệ thống được mô tả bằng các thành phần thực thi chính và các kiểu liên kết giữa chúng. Góc nhìn này đặc biệt quan trọng vì nó cho biết một yêu cầu đi qua những khối nào, khối nào giữ quyết định nghiệp vụ và khối nào chỉ đóng vai trò tích hợp hoặc hỗ trợ kỹ thuật.

\subsubsection{Các thành phần chính}\label{sec:main-components}
Góc nhìn này bóc tách hệ thống thành các khối thực thi cụ thể, đảm bảo sự phân tách rạch ròi giữa trải nghiệm giao diện, điều phối luồng, luật nghiệp vụ và tích hợp hạ tầng phần cứng lẫn phần mềm:

\begin{itemize}
  \item \textbf{Các ứng dụng giao diện (Client Applications):} Gồm giao diện cho lễ tân, phục vụ, bếp, thu ngân và quản lý. Lớp này tiếp nhận thao tác người dùng, gọi các endpoint trong \textsf{ENDPOINTS} của frontend và hiển thị lại dữ liệu do backend trả về:
  \begin{itemize}
      \item Dùng \textsf{apiFetch} và \textsf{ENDPOINTS} để gọi REST API qua \textsf{Nitro API Proxy}.
      \item Dùng \textsf{permissions.ts} và dữ liệu vai trò từ \textsf{/api/auth/me} để kiểm soát điều hướng theo quyền trên giao diện.
  \end{itemize}

  \item \textbf{Cổng vào hệ thống (API Gateway \& Frontend Proxy):} \textsf{Frontend Node Server} phục vụ React và chuyển tiếp \textsf{/api}; \textsf{GatewayProxyController} trong \textsf{api-gateway} định tuyến yêu cầu REST đến các service phía sau. Thiết kế hiện tại chưa cho thấy WebSocket/SSE, rate limiting hay cấu hình Nginx riêng.

  \item \textbf{Các dịch vụ ứng dụng (Application Services):} Lớp này chịu trách nhiệm điều phối các use case. Thành phần này không chứa luật nghiệp vụ mà chỉ quản lý ranh giới giao dịch, gọi các đối tượng miền phù hợp và ủy quyền việc phát sự kiện.

  \item \textbf{Các service miền nghiệp vụ (Domain Services):} Đây là lớp lõi của kiến trúc, chứa các Aggregate Roots và các bất biến nghiệp vụ. Mọi quyết định nghiệp vụ như tính tiền hoặc xác định trạng thái bàn chỉ được xử lý tại lớp này.

  \item \textbf{Lớp lưu trữ dữ liệu (Data Persistence):} Thực hiện lưu trữ thông qua mẫu Repository nhằm cô lập nghiệp vụ khỏi nền tảng SQL. Lớp này phân tách rõ giữa \textbf{Mô hình ghi (Transactional Schema)} phục vụ cho tốc độ giao dịch nóng và \textbf{Mô hình đọc (Reporting Projections)} phục vụ truy vấn phân tích để không gây nghẽn cổ chai cơ sở dữ liệu.

  \item \textbf{Bộ truyền sự kiện và xử lý nền (Event Bus \& Background Workers):} Đảm bảo tính nhất quán cuối cùng (Eventual Consistency) của hệ thống. Hệ thống áp dụng \textbf{Outbox Pattern} với \textbf{Outbox Dispatcher} để quét và đẩy sự kiện một cách đáng tin cậy. Các hệ quả bất đồng bộ (như cập nhật kho, gửi thông báo) được đưa vào hàng đợi nền để không làm chậm luồng thao tác trực tiếp của nhân viên.

  \item \textbf{Các bộ kết nối thích nghi (Adapters):} Được thiết kế theo nguyên lý Đảo ngược phụ thuộc (Dependency Inversion), chia làm hai nhóm đã thấy trong thiết kế nhằm cách ly lõi nghiệp vụ:
  \begin{itemize}
      \item \textit{Bộ kết nối thanh toán}: \textsf{CashPaymentGateway} và \textsf{ExternalReferencePaymentGateway} hiện thực interface \textsf{PaymentGateway}.
      \item \textit{Bộ kết nối thông báo và biên lai}: \textsf{EmailNotificationChannelAdapter}, \textsf{SmsNotificationChannelAdapter}, \textsf{InAppNotificationChannelAdapter}, \textsf{PrintNotificationChannelAdapter}, \textsf{PrintReceiptDeliveryAdapter} và \textsf{DigitalReceiptDeliveryAdapter}.
  \end{itemize}
\end{itemize}

\subsubsection{Các kiểu kết nối chính}\label{sec:main-connectors}
Để đáp ứng yêu cầu về hiệu năng cao và tính nhất quán của dữ liệu trong môi trường vận hành nhà hàng, IRMS kết hợp các kiểu kết nối đã có trong thiết kế và cấu hình. Các tương tác giữa các thành phần được định nghĩa qua bảng sau:

\begin{table}[H]
\centering
\caption{Các kiểu kết nối chính trong IRMS}
\begin{tabularx}{\textwidth}{|p{2.95cm}|X|p{4cm}|}
\hline
\textbf{Kiểu kết nối} & \textbf{Vai trò và Đặc tính kỹ thuật} & \textbf{Ví dụ thực tế trong IRMS} \\
\hline
Giao thức HTTP (REST API) & Kết nối đồng bộ (Synchronous), ưu tiên độ trễ thấp và phản hồi trạng thái giao dịch tức thì. & POS App gửi yêu cầu xác nhận đơn gọi món hoặc lập hóa đơn. \\
\hline
REST qua \textsf{apiFetch} & Frontend gọi định kỳ hoặc gọi lại sau thao tác người dùng để tải trạng thái mới. & Màn hình bếp gọi \textsf{/api/kitchen/overview}; màn hình báo cáo gọi các endpoint trong \textsf{/api/reports}. \\
\hline
Thông điệp Sự kiện (Pub/Sub) & Kết nối bất đồng bộ (Asynchronous) qua Message Broker, giúp giảm sự phụ thuộc trực tiếp (Decoupling) giữa các module. & Lan truyền sự kiện đơn hàng hoàn tất để trừ tồn kho và làm giàu dữ liệu báo cáo. \\
\hline  
Adapter nội bộ & Interface và adapter trong thiết kế để cô lập biến thể nghiệp vụ hoặc kênh phát hành. & \textsf{PaymentGateway}, \textsf{NotificationChannelAdapter} và \textsf{ReceiptDeliveryAdapter}. \\
\hline
Kết nối Database (JDBC / Driver) & Truy xuất dữ liệu có cấu trúc và quản lý ranh giới giao dịch (Atomic Transactions). & Business Services lưu trữ trạng thái bàn và phiên phục vụ vào PostgreSQL. \\
\hline
\end{tabularx}
\end{table}

\subsubsection{Cơ sở lựa chọn mô hình kết nối hỗn hợp}\label{sec:connector-rationale}
Kiến trúc IRMS không sử dụng duy nhất một kiểu kết nối mà kết hợp linh hoạt dựa trên hai nguyên tắc cốt lõi:

\begin{enumerate}
    \item \textbf{Tuyến đầu ưu tiên đồng bộ (Sync for UX):} Các thao tác tương tác trực tiếp với khách hàng (đặt bàn, gọi món, thanh toán) được thực hiện qua API đồng bộ. Điều này giúp nhân viên nhận được phản hồi "Thành công" hoặc "Thất bại" ngay lập tức, đảm bảo luồng phục vụ không bị gián đoạn.
    
    \item \textbf{Hậu phương ưu tiên bất đồng bộ (Async for Resilience):} Các hệ quả phát sinh sau giao dịch như gửi thông báo, cập nhật kho, audit và tổng hợp báo cáo được đẩy sang tuyến sự kiện qua \textsf{RabbitMQ}, \textsf{outbox\_events} và các bộ xử lý nền. Việc này giúp bảo vệ hệ thống: nếu một tác vụ phụ bị chậm hoặc lỗi, nó cũng không làm nghẽn luồng nghiệp vụ chính.
\end{enumerate}


\textbf{Góc nhìn chuyên sâu về thành phần và kết nối:} \\
Điểm then chốt trong kiến trúc IRMS là sự phân định rạch ròi giữa \emph{kết nối phục vụ trải nghiệm người dùng} và \emph{kết nối phục vụ phối hợp nội bộ}. Việc chọn đúng chiến lược giao tiếp cho từng ngữ cảnh không chỉ quyết định độ trễ phản hồi mà còn ảnh hưởng trực tiếp đến khả năng chịu lỗi (fault tolerance) và tính mở rộng của toàn hệ thống.



\begin{figure}[H]
\centering
\IfFileExists{09_component_and_connector_view_overall_runtime_structure.pdf}{
  \includegraphics[width=0.97\linewidth,height=0.84\textheight,keepaspectratio]{09_component_and_connector_view_overall_runtime_structure.pdf}
}{
  \IfFileExists{09_component_and_connector_view_overall_runtime_structure.png}{
    \includegraphics[width=0.97\linewidth,height=0.84\textheight,keepaspectratio]{09_component_and_connector_view_overall_runtime_structure.png}
  }{
    \fbox{\parbox{0.82\linewidth}{Thiếu file \textsf{09\_component\_and\_connector\_view\_overall\_runtime\_structure.(pdf/png)}.}}
  }
}
\caption{Sơ đồ thành phần và kết nối tổng quan của IRMS}
\label{fig:irms-cc-overall}
\end{figure}



Hình~\ref{fig:irms-cc-overall} thể hiện cấu trúc thực thi tổng quan của hệ thống theo cách nhìn thành phần. Thiết bị người dùng chỉ giao tiếp với lớp cổng vào; từ đó yêu cầu đi qua lớp điều phối ca sử dụng, xuống miền nghiệp vụ, lớp lưu trữ và các adapter biên. Ý nghĩa học thuật của hình này là làm rõ hệ thống không được tổ chức như một khối gọi chéo tùy ý, mà có hướng phụ thuộc rõ ràng từ ngoài vào trong.

\paragraph{Các thành phần chính}
\begin{itemize}
  \item \textbf{Các ứng dụng giao diện}: là nơi phát sinh yêu cầu từ lễ tân, phục vụ, bếp, thu ngân và quản lý.
  \item \textbf{Cổng vào giao diện lập trình ứng dụng}: tiếp nhận yêu cầu, xác thực, kiểm tra đầu vào và chuyển tiếp tới đúng ca sử dụng.
  \item \textbf{Các dịch vụ ứng dụng}: điều phối từng ca sử dụng cụ thể như tạo đặt bàn, xác nhận đơn gọi món, cập nhật phiếu bếp hoặc chốt hóa đơn.
  \item \textbf{Các service miền nghiệp vụ}: giữ luật nghiệp vụ cốt lõi, bất biến dữ liệu và quyết định nghiệp vụ của nhà hàng.
  \item \textbf{Lớp lưu trữ dữ liệu}: chịu trách nhiệm lưu trữ trạng thái giao dịch và bảo vệ ranh giới giao dịch.
  \item \textbf{Bộ phát sự kiện và bộ xử lý nền}: xử lý sự kiện và tác vụ nền không nên nằm trên tuyến giao dịch nóng.
  \item \textbf{Các port/adapter biên}: bao bọc cách xử lý thanh toán, gửi thông báo và phát hành biên lai.
\end{itemize}


\begin{figure}[H]
\centering
\IfFileExists{10_component_diagram_detailed_backend_composition.pdf}{
  \includegraphics[width=0.97\linewidth,height=0.84\textheight,keepaspectratio]{10_component_diagram_detailed_backend_composition.pdf}
}{
  \IfFileExists{10_component_diagram_detailed_backend_composition.png}{
    \includegraphics[width=0.97\linewidth,height=0.84\textheight,keepaspectratio]{10_component_diagram_detailed_backend_composition.png}
  }{
    \fbox{\parbox{0.82\linewidth}{Thiếu file \textsf{10\_component\_diagram\_detailed\_backend\_composition.(pdf/png)}.}}
  }
}
\caption{Sơ đồ thành phần chi tiết của lớp service nghiệp vụ trong IRMS}
\label{fig:irms-cc-backend}
\end{figure}



Hình~\ref{fig:irms-cc-backend} mở chi tiết lớp backend ở mức component, nhưng không trộn sang Class Diagram. Sơ đồ gom các controller REST, application service, required interface và provided adapter theo đúng vai trò trong thiết kế. Nhờ đó người đọc thấy rõ yêu cầu đi từ biên API vào lớp điều phối, sau đó phụ thuộc vào các port như repository, \textsf{PaymentGateway}, \textsf{ReceiptDeliveryAdapter}, \textsf{SharedIdentity*Port} và \textsf{DomainEventPublisher}.

Điểm chính của hình là hướng phụ thuộc được kiểm soát qua interface. Các application service không phụ thuộc trực tiếp vào JDBC adapter, gateway thanh toán hay cơ chế outbox cụ thể; những thành phần đó cung cấp lại năng lực qua adapter ở rìa. Cách biểu diễn này khớp với thiết kế hiện tại, nơi các lớp như \textsf{JdbcOrderRepository}, \textsf{CashPaymentGateway}, \textsf{ExternalReferencePaymentGateway}, \textsf{RemoteSharedIdentity*Client} và \textsf{TransactionalOutboxPublisher} nằm ở phía hiện thực kỹ thuật.

\paragraph{Các thành phần chính}
\begin{itemize}
  \item \textbf{API layer}: gom các controller REST của những module nghiệp vụ chính.
  \item \textbf{Application services}: gom các service điều phối use case như \textsf{OrdersService}, \textsf{KitchenService}, \textsf{BillingService}, \textsf{ReservationService}, \textsf{InventoryService}, \textsf{ReportingQueryService}, \textsf{AuthService} và \textsf{AuditService}.
  \item \textbf{Required interfaces}: mô tả các port mà lớp ứng dụng cần, gồm port truy cập dữ liệu, integration port và event publisher.
  \item \textbf{Provided adapters}: mô tả các adapter cung cấp các port trên, gồm JDBC adapter, payment/receipt adapter, shared identity client và outbox publisher.
\end{itemize}



\subsubsection{Sơ đồ thành phần cho tuyến lệnh vận hành}


\begin{figure}[H]
\centering
\IfFileExists{24_component_diagram_operational_command_collaboration.pdf}{
  \includegraphics[width=0.97\linewidth,height=0.84\textheight,keepaspectratio]{24_component_diagram_operational_command_collaboration.pdf}
}{
  \IfFileExists{24_component_diagram_operational_command_collaboration.png}{
    \includegraphics[width=0.97\linewidth,height=0.84\textheight,keepaspectratio]{24_component_diagram_operational_command_collaboration.png}
  }{
    \fbox{\parbox{0.8\linewidth}{Thiếu file \textsf{24\_component\_diagram\_operational\_command\_collaboration.(pdf/png)}.}}
  }
}
\caption{Sơ đồ thành phần cho tuyến lệnh vận hành của IRMS}
\label{fig:irms-component-operational-command}
\end{figure}



Hình~\ref{fig:irms-component-operational-command} mô tả luồng thực thi của một lệnh vận hành (operational command), từ giao diện người dùng đến khi dữ liệu được lưu trữ và các tác vụ phụ (side-effects) được kích hoạt. Trình tự này được thiết kế chặt chẽ để đảm bảo hệ thống hoạt động ổn định và có khả năng truy vết (audit).

Đầu tiên, các yêu cầu đi qua bộ lọc xác thực và phân quyền ở tầng biên (edge) để kiểm tra tính hợp lệ trước khi chạm vào hệ thống lõi. Sau đó, request được chuyển đến các Application Service (\textsf{Reservation}, \textsf{Ordering}, \textsf{Kitchen}, \textsf{Billing}). Tại đây, Application Service chỉ đóng vai trò điều phối (orchestration): gọi các đối tượng Domain tương ứng, quản lý vòng đời giao dịch (transaction), và ghi sự kiện vào Outbox (Outbox Pattern) để chuẩn bị cho các tác vụ tiếp theo.

Thứ hai, thiết kế tách biệt rõ ràng giữa \textbf{giao dịch lõi} (core transaction) và \textbf{tác vụ phụ} (side-effects). Các giao dịch chính (như chốt đơn, ghi nhận thanh toán) bắt buộc phải được commit thành công vào cơ sở dữ liệu. Chỉ khi trạng thái lõi đã nhất quán, hệ thống mới phát sự kiện (event) để xử lý các tác vụ bất đồng bộ như báo xuống bếp, gửi thông báo hoặc cập nhật Read Model. Cơ chế này thực thi nguyên tắc \enquote{lõi đồng bộ, hỗ trợ bất đồng bộ}, giúp tối ưu thời gian phản hồi cho client nhưng vẫn đảm bảo toàn vẹn dữ liệu.

Cuối cùng, việc đặt \textsf{PaymentGateway}, \textsf{ReceiptDeliveryAdapter} và \textsf{NotificationChannelAdapter} ở rìa kiến trúc thể hiện rõ việc áp dụng mô hình Ports \& Adapters. Tầng Domain chỉ định nghĩa các interface (Port), trong khi chi tiết tích hợp cụ thể với nhà cung cấp thanh toán hay máy in biên lai được giao cho các Adapter bên ngoài. Cách tổ chức này giúp lõi nghiệp vụ độc lập hoàn toàn với chi tiết hạ tầng.
\paragraph{Các thành phần chính}
\begin{itemize}
  \item \textbf{Cổng vào giao diện lập trình ứng dụng và lớp bộ điều khiển} cùng \textbf{bộ lọc xác thực / kiểm tra chính sách}: tiếp nhận yêu cầu, chuẩn hóa dữ liệu vào, xác thực và chặn hành vi vượt quyền trước khi đi vào lõi.
  \item \textbf{Các dịch vụ ứng dụng} của từng vùng \textsf{Reservation}, \textsf{Ordering}, \textsf{Kitchen} và \textsf{Billing}: điều phối use case, gọi đúng thành phần miền nghiệp vụ và kiểm soát biên giao dịch.
  \item \textbf{Các thành phần miền nghiệp vụ} và \textbf{kho lưu trữ}: hiện thực quy tắc miền và nơi lưu trạng thái bền vững.
  \item \textbf{Hộp thư ra / bộ phát sự kiện} cùng các port/adapter như \textsf{PaymentGateway}, \textsf{ReceiptDeliveryAdapter} và \textsf{NotificationChannelAdapter}: tách bước chốt giao dịch khỏi bước xử lý biên hoặc xử lý nền.
\end{itemize}


\subsubsection{Sơ đồ thành phần cho xử lý nền, quản trị và mô hình đọc}
\begin{figure}[H]
\centering
\IfFileExists{25_component_diagram_supporting_async_and_governance.pdf}{
  \includegraphics[width=0.97\linewidth,height=0.84\textheight,keepaspectratio]{25_component_diagram_supporting_async_and_governance.pdf}
}{
  \IfFileExists{25_component_diagram_supporting_async_and_governance.png}{
    \includegraphics[width=0.97\linewidth,height=0.84\textheight,keepaspectratio]{25_component_diagram_supporting_async_and_governance.png}
  }{
    \fbox{\parbox{0.82\linewidth}{Thiếu file \textsf{25\_component\_diagram\_supporting\_async\_and\_governance.(pdf/png)}.}}
  }
}
\caption{Sơ đồ thành phần cho xử lý nền, quản trị và mô hình đọc}
\label{fig:irms-component-supporting-async}
\end{figure}


Hình~\ref{fig:irms-component-supporting-async} mô tả kiến trúc event-driven sau khi các giao dịch lõi đã được xác lập. Đây là Component Diagram của C\&C View, không phải Class Diagram: các phần tử trong hình là service, outbox, RabbitMQ và bộ xử lý sự kiện nền đang tham gia vào tương tác runtime. Tại đây, IRMS tập trung vào việc hiện thực hóa nguyên tắc \textbf{Nhất quán cuối cùng (Eventual Consistency)} và \textbf{Tách biệt trách nhiệm truy vấn (CQRS - Command Query Responsibility Segregation)} ở mức độ thực dụng.

Cấu trúc phối hợp giữa các thành phần hỗ trợ được tổ chức như sau:

\begin{enumerate}
    \item \textbf{Cơ chế phát và chuyển tiếp sự kiện:} Các service giao dịch ghi sự kiện vào \textsf{outbox\_events}. Sau đó \textsf{RabbitMqOutboxRelay} đọc bản ghi chưa phát, dùng \textsf{RabbitMqEventPublisher} đẩy sự kiện sang \textsf{RabbitMQ}. Cấu trúc này khớp với \texttt{TransactionalOutboxPublisher}, \texttt{RabbitMqOutboxRelay} và cấu hình \texttt{spring.rabbitmq.*} trong backend.

    \item \textbf{Cơ chế diễn dịch sự kiện (Event Projection):} Các sự kiện sau khi đi qua \textsf{RabbitMQ} được tiêu thụ bởi các bộ xử lý sự kiện như \textsf{ReportingProjectionConsumer} và \textsf{InventoryStockChangedConsumer}. Các thành phần này chuyển dữ liệu từ dạng giao dịch sang mô hình đọc hoặc projection phục vụ báo cáo và tồn kho, giúp yêu cầu từ màn hình quản lý không tranh chấp trực tiếp với luồng gọi món.
    
    \item \textbf{Quản trị và Tuân thủ (Governance \& Compliance):} 
    \begin{itemize}
        \item \textsf{AuditRecordingRequestedConsumer}: Độc lập với logic nghiệp vụ, chỉ tiêu thụ event để ghi audit log (lưu vết: ai, làm gì, khi nào, ngữ
        \item \textbf{IAM / Authorization Service:} Cung cấp các chính sách bảo vệ tài nguyên (Policy-based Access Control), đảm bảo các thao tác quản trị nhạy cảm như hoàn tiền hoặc điều chỉnh tồn kho luôn được kiểm soát chặt chẽ.
    \end{itemize}

    \item \textbf{Tích hợp thông báo và liên lạc (Notification Hub):} \textsf{NotificationRequestedConsumer} tiếp nhận sự kiện từ \textsf{RabbitMQ}; sau đó \textsf{NotificationService} ghi nhận thông báo qua các adapter kênh trong thiết kế:
    \begin{itemize}
        \item \textsf{EmailNotificationChannelAdapter}, \textsf{SmsNotificationChannelAdapter};
        \item \textsf{InAppNotificationChannelAdapter}, \textsf{PrintNotificationChannelAdapter}.
    \end{itemize}
    Thiết kế hiện tại chưa thể hiện nhà cung cấp SMS/Email bên ngoài hay máy in vật lý, nên phần này chỉ được hiểu là biên adapter nội bộ.
\end{enumerate}

\paragraph{Cơ sở thiết kế xử lý nền:}
\begin{itemize}
  \item \textbf{Tối ưu hóa tài nguyên ghi (Write Optimization):} Đẩy các tác vụ phụ (tổng hợp báo cáo, gửi thông báo) sang xử lý nền giúp giảm tải DB giao dịch chính, tối ưu tốc độ phản hồi cho các nghiệp vụ lõi (chốt đơn, thanh toán).
  \item \textbf{Khả năng quan sát (Observability):} Tách biệt các event consumer giúp dễ dàng monitor trạng thái và độ trễ của từng tác vụ (cập nhật projection, audit log, notification) thông qua log và các bảng \textsf{outbox\_events}, \textsf{processed\_events}.
\end{itemize}

\subsection{Góc nhìn triển khai}\label{sec:deployment-view}
\subsubsection{Tổng quan phân bổ hệ thống}

Ở góc nhìn triển khai, hệ thống được ánh xạ lên hạ tầng vật lý và hạ tầng thực thi nhằm làm rõ cách các thành phần phần mềm được phân bổ trên các nút triển khai (deployment nodes), cách các nút này giao tiếp với nhau và cách hệ thống đáp ứng các yêu cầu phi chức năng như độ trễ, khả năng mở rộng và độ tin cậy. Mô hình này nhất quán với \adrref{ADR-0008}{adr:0008-single-postgresql} về một PostgreSQL vật lý duy nhất và \adrref{ADR-0014}{adr:0014-docker-compose} về Docker Compose cho runtime cục bộ.

Hệ thống được tổ chức thành bốn tầng chính như sau:

\begin{enumerate}

  \item \textbf{Tầng thiết bị đầu cuối (Client Layer)}
  \begin{itemize}
    \item Máy tính bảng dành cho nhân viên phục vụ, hỗ trợ thao tác đặt bàn và gọi món theo thời gian thực.
    \item Thiết bị lễ tân hoặc ki-ốt tiếp nhận khách, phục vụ việc đăng ký và phân bổ chỗ ngồi.
    \item Màn hình bếp tại từng trạm chế biến, hiển thị danh sách món cần chuẩn bị với yêu cầu cập nhật gần thời gian thực.
    \item Máy thu ngân, thực hiện các giao dịch thanh toán với yêu cầu độ trễ thấp và tính nhất quán cao.
    \item Trình duyệt của quản lý, phục vụ truy vấn báo cáo và giám sát hoạt động hệ thống.
  \end{itemize}

  Các thiết bị này hoạt động như các client gửi yêu cầu đồng bộ (synchronous request) tới hệ thống thông qua giao thức HTTP/HTTPS. Đặc điểm chung của tầng này là phụ thuộc vào mạng và yêu cầu độ trễ thấp đối với các thao tác nghiệp vụ chính.

  \item \textbf{Tầng cổng vào và phân phối (Gateway \& Delivery Layer)}
  \begin{itemize}
    \item Bộ đảo ngược lưu lượng hoặc \textbf{API Gateway}, đóng vai trò điểm truy cập duy nhất của hệ thống.
    \item Máy chủ phân phối giao diện tĩnh hoặc giao diện phía máy chủ, chịu trách nhiệm phục vụ nội dung frontend.
  \end{itemize}

  Tầng này chịu trách nhiệm:
  \begin{itemize}
    \item Điều phối yêu cầu từ client tới các service tương ứng.
    \item Thực hiện các kiểm tra kỹ thuật chung như xác thực, phân quyền và ghi nhận nhật ký truy cập.
    \item Giảm tải cho các service phía sau thông qua việc gom và chuẩn hóa các điểm truy cập.
  \end{itemize}

  Việc sử dụng API Gateway giúp cô lập tầng client với tầng ứng dụng, đồng thời tạo điều kiện thuận lợi cho việc mở rộng và thay đổi cấu trúc backend trong tương lai.

  \item \textbf{Tầng ứng dụng (Application Layer)}
  \begin{itemize}
    \item Khối xử lý phía sau bao gồm các service nghiệp vụ chính:
    \begin{itemize}
      \item Service đặt bàn và quản lý chỗ ngồi.
      \item Service gọi món và thực đơn.
      \item Service điều phối bếp.
      \item Service thanh toán.
      \item Service quản lý tồn kho.
      \item Service báo cáo và quản trị.
      \item Service quản lý truy cập và kiểm toán.
    \end{itemize}

    \item Bộ xử lý nền (Background Worker) thực hiện các tác vụ bất đồng bộ:
    \begin{itemize}
      \item Xử lý sự kiện phát sinh từ hệ thống (event-driven processing).
      \item Thực hiện các tác vụ retry khi xảy ra lỗi tạm thời.
      \item Làm mới các ảnh chụp dữ liệu phục vụ báo cáo (reporting projections).
    \end{itemize}
  \end{itemize}

  Các service trong tầng này có thể được triển khai trên cùng một cụm máy chủ hoặc container, nhưng được tách biệt về mặt logic theo miền nghiệp vụ. Các giao tiếp giữa service có thể bao gồm:
  \begin{itemize}
    \item Giao tiếp đồng bộ (REST API) cho các thao tác nghiệp vụ chính.
    \item Giao tiếp bất đồng bộ (message queue) cho các tác vụ nền.
  \end{itemize}

  \item \textbf{Tầng dữ liệu (Data Layer)}
  \begin{itemize}
    \item Cơ sở dữ liệu giao dịch, lưu trữ dữ liệu vận hành chính như đơn hàng, thanh toán, tồn kho và người dùng.
    \item Mô hình đọc và khung nhìn báo cáo, được tối ưu cho truy vấn và phân tích dữ liệu.
    \item Các bảng \textsf{outbox\_events}, \textsf{processed\_events} và mô hình đọc báo cáo phục vụ xử lý nền và truy vấn quản trị.
    \item Hạ tầng \textsf{RabbitMQ} phục vụ hàng đợi sự kiện giữa các service.
  \end{itemize}

  Tầng dữ liệu được thiết kế tách biệt với tầng ứng dụng nhằm đảm bảo tính toàn vẹn dữ liệu, đồng thời hỗ trợ mở rộng theo nhu cầu đọc/ghi. Các thành phần như mô hình đọc và hệ thống logging thường được cập nhật thông qua luồng xử lý bất đồng bộ để giảm tải cho hệ thống giao dịch chính.

\end{enumerate}

\subsubsection{Cơ sở thiết kế triển khai}
\begin{itemize}
  \item \textbf{Tối ưu mạng nội bộ:} Thiết bị gọi món và bếp ưu tiên kết nối mạng LAN để giảm latency và đảm bảo vận hành liên tục khi kết nối WAN (Internet) không ổn định.
  \item \textbf{Độ sẵn sàng cao (High Availability):} Các module Billing và Ordering yêu cầu HA cao nhất. Cần ưu tiên giám sát, backup và khả năng scale vì ảnh hưởng trực tiếp đến doanh thu.
  \item \textbf{Tách biệt OLTP và OLAP:} Tách khối báo cáo khỏi tuyến giao dịch để các truy vấn phân tích (nặng) không gây nghẽn luồng xử lý gọi món và thanh toán (thời gian thực).
  \item \textbf{Bảo mật dữ liệu:} Áp dụng chính sách phân quyền và sao lưu riêng biệt cho dữ liệu thanh toán và audit log; đây là nhóm dữ liệu nhạy cảm, yêu cầu bảo vệ cao nhất.
  \item \textbf{Cô lập tài nguyên (Resource Isolation):} Tách bộ xử lý nền (Background Worker) giúp các tác vụ phụ như in biên lai điện tử hay gửi thông báo không tranh chấp tài nguyên với các thao tác tại quầy.
\end{itemize}



\begin{figure}[H]
\centering
\IfFileExists{15_deployment_diagram_physical_deployment_overview.pdf}{
  \includegraphics[width=0.97\linewidth,height=0.84\textheight,keepaspectratio]{15_deployment_diagram_physical_deployment_overview.pdf}
}{
  \IfFileExists{15_deployment_diagram_physical_deployment_overview.png}{
    \includegraphics[width=0.97\linewidth,height=0.84\textheight,keepaspectratio]{15_deployment_diagram_physical_deployment_overview.png}
  }{
    \fbox{\parbox{0.82\linewidth}{Thiếu file \textsf{15\_deployment\_diagram\_physical\_deployment\_overview.(pdf/png)}.}}
  }
}
\caption{Sơ đồ triển khai tổng quan về bố trí vật lý của IRMS}
\label{fig:irms-deploy-overall}
\end{figure}

Hình~\ref{fig:irms-deploy-overall} mô tả cách hệ thống được triển khai trong môi trường thực tế của nhà hàng. 

Điểm quan trọng nhất cần rút ra là IRMS được tổ chức theo mô hình \textbf{nhiều điểm chạm vật lý nhưng một lõi điều phối tập trung}. Các thiết bị tại hiện trường --- như máy tính bảng phục vụ, thiết bị lễ tân, màn hình bếp và quầy thu ngân --- được đặt gần người dùng để đảm bảo thao tác nhanh và thuận tiện. Tuy nhiên, toàn bộ trạng thái nghiệp vụ cốt lõi vẫn được duy trì tại khối xử lý trung tâm và hạ tầng dữ liệu, thay vì phân tán xuống các thiết bị.

Cách bố trí này giúp tránh tình trạng mất nhất quán khi mạng dao động hoặc khi nhiều thao tác xảy ra đồng thời tại các điểm khác nhau. Nó cũng phản ánh một nguyên tắc quan trọng: các thiết bị đầu cuối chỉ nên đóng vai trò tương tác, còn quyết định nghiệp vụ bền vững phải được neo tại backend.

Một đặc điểm thiết kế quan trọng khác là sự tách biệt giữa hai tuyến xử lý:

\begin{itemize}
  \item \textbf{Tuyến xử lý trực tuyến (synchronous path)} phục vụ các thao tác yêu cầu phản hồi ngay như gọi món, cập nhật trạng thái bếp và thanh toán.
  \item \textbf{Tuyến xử lý bất đồng bộ (asynchronous path)} xử lý các hệ quả phụ thông qua \textbf{Outbox / hàng đợi} và bộ xử lý nền, bao gồm gửi thông báo, phát hành biên lai và cập nhật dữ liệu báo cáo.
\end{itemize}

Việc phân tách này cho phép hệ thống giữ được thời gian phản hồi ổn định trong giờ cao điểm, khi các thao tác trực tuyến cần được ưu tiên tuyệt đối.

Ngoài ra, sơ đồ cũng làm rõ vai trò của \textsf{PostgreSQL}, \textsf{RabbitMQ}, \textsf{outbox\_events} và các bảng hoặc view báo cáo. Các tác vụ không yêu cầu phản hồi ngay được xử lý tách biệt, trong khi dữ liệu giao dịch chính vẫn được neo trong PostgreSQL để giữ tính nhất quán.

Từ góc độ vận hành, cách bố trí này còn mang lại khả năng \textbf{cô lập lỗi theo thành phần}. Khi bộ xử lý sự kiện nền, hàng đợi hoặc mô hình đọc báo cáo gặp sự cố, tuyến xử lý trực tuyến vẫn có thể tiếp tục hoạt động với mức suy giảm có kiểm soát.

\paragraph{Các thành phần chính}
\begin{itemize}
  \item \textbf{Thiết bị đầu cuối}: máy tính bảng phục vụ, thiết bị lễ tân, màn hình bếp, máy thu ngân và trình duyệt quản lý.
  \item \textbf{Khối cổng vào}: \textsf{frontend container} và \textsf{api-gateway container} tiếp nhận, proxy và định tuyến yêu cầu từ thiết bị.
  \item \textbf{Khối xử lý trung tâm}: các container service Spring Boot thực hiện logic nghiệp vụ chính.
  \item \textbf{Hạ tầng sự kiện}: \textsf{RabbitMQ}, \textsf{outbox\_events} và \textsf{processed\_events} xử lý các tác vụ bất đồng bộ.
  \item \textbf{Hạ tầng dữ liệu}: \textsf{PostgreSQL} lưu dữ liệu giao dịch và mô hình đọc báo cáo.
\end{itemize}

\paragraph{Lý do thiết kế.}
\begin{itemize}
  \item Giữ một lõi xử lý trung tâm giúp đảm bảo tính nhất quán và đơn giản hóa triển khai.
  \item Tách tuyến bất đồng bộ giúp giảm tải cho tuyến trực tuyến.
  \item Phân lớp lưu trữ giúp hệ thống linh hoạt và dễ mở rộng hơn.
\end{itemize}


% =========================

\begin{figure}[H]
\centering
\IfFileExists{16_deployment_diagram_network_zones_and_trust_boundaries.pdf}{
  \includegraphics[width=0.97\linewidth,height=0.84\textheight,keepaspectratio]{16_deployment_diagram_network_zones_and_trust_boundaries.pdf}
}{
  \IfFileExists{16_deployment_diagram_network_zones_and_trust_boundaries.png}{
    \includegraphics[width=0.97\linewidth,height=0.84\textheight,keepaspectratio]{16_deployment_diagram_network_zones_and_trust_boundaries.png}
  }{
    \fbox{\parbox{0.82\linewidth}{Thiếu file \textsf{16\_deployment\_diagram\_network\_zones\_and\_trust\_boundaries.(pdf/png)}.}}
  }
}
\caption{Sơ đồ triển khai theo vùng mạng và ranh giới tin cậy}
\label{fig:irms-deploy-network}
\end{figure}

Hình~\ref{fig:irms-deploy-network} mở rộng góc nhìn triển khai bằng cách làm rõ \textbf{ranh giới tin cậy} và cách hệ thống được phân chia thành các vùng mạng khác nhau.

Kiến trúc được chia thành các vùng: vùng thiết bị hiện trường, vùng ứng dụng, vùng dữ liệu và vùng hàng đợi sự kiện. Mỗi vùng đại diện cho một mức độ tin cậy và quyền truy cập khác nhau.

Nguyên tắc quan trọng nhất là \textbf{không có thành phần nào được phép vượt qua ranh giới mà không qua kiểm soát}. Các thiết bị đầu cuối chỉ giao tiếp với hệ thống thông qua \textsf{frontend container} và \textsf{api-gateway}. Các service phía sau mới được truy cập \textsf{PostgreSQL} và \textsf{RabbitMQ}.

Việc tách riêng vùng dữ liệu giúp bảo vệ dữ liệu nhạy cảm và giảm nguy cơ truy cập trái phép. Đồng thời, việc đặt \textsf{RabbitMQ} ở vùng riêng giúp cô lập tuyến sự kiện khỏi tuyến yêu cầu trực tuyến.

Sơ đồ này cũng hỗ trợ vận hành thực tế: khi xảy ra sự cố, việc phân vùng giúp nhanh chóng xác định lỗi thuộc lớp nào.

\paragraph{Các vùng chính.}
\begin{itemize}
  \item Vùng thiết bị hiện trường
  \item Vùng ứng dụng
  \item Vùng dữ liệu
  \item Vùng dữ liệu và hàng đợi sự kiện
\end{itemize}

\paragraph{Lý do thiết kế.}
\begin{itemize}
  \item Tăng cường bảo mật thông qua phân tách vùng
  \item Kiểm soát luồng truy cập giữa các thành phần
  \item Hỗ trợ vận hành và xử lý sự cố
\end{itemize}

% =========================

\begin{figure}[H]
\centering
\IfFileExists{17_deployment_diagram_data_and_execution_mapping.pdf}{
  \includegraphics[width=0.97\linewidth,height=0.84\textheight,keepaspectratio]{17_deployment_diagram_data_and_execution_mapping.pdf}
}{
  \IfFileExists{17_deployment_diagram_data_and_execution_mapping.png}{
    \includegraphics[width=0.97\linewidth,height=0.84\textheight,keepaspectratio]{17_deployment_diagram_data_and_execution_mapping.png}
  }{
    \fbox{\parbox{0.82\linewidth}{Thiếu file \textsf{17\_deployment\_diagram\_data\_and\_execution\_mapping.(pdf/png)}.}}
  }
}
\caption{Sơ đồ triển khai về ánh xạ dữ liệu và thực thi}
\label{fig:irms-deploy-data}
\end{figure}

Hình~\ref{fig:irms-deploy-data} làm rõ cách các service, dữ liệu và luồng thực thi được tổ chức trong hệ thống.

Sơ đồ này thể hiện rõ rằng IRMS theo \textbf{kiến trúc hướng dịch vụ kết hợp điều phối bằng sự kiện}. Các service nghiệp vụ xử lý giao dịch chính thông qua API đồng bộ, trong khi các hệ quả phụ được xử lý thông qua cơ chế bất đồng bộ dựa trên \textbf{Outbox và hàng đợi sự kiện}.

Luồng xử lý có thể chia thành hai phần:

\begin{itemize}
  \item \textbf{Luồng giao dịch chính}: các service ghi dữ liệu vào cơ sở dữ liệu giao dịch.
  \item \textbf{Luồng bất đồng bộ}: các sự kiện được ghi vào Outbox, sau đó được bộ xử lý nền tiêu thụ để cập nhật mô hình đọc, gửi thông báo hoặc xử lý các tác vụ phụ.
\end{itemize}

Cách tổ chức này giúp đảm bảo rằng các giao dịch cốt lõi vẫn rõ ràng và nhất quán, trong khi hệ thống vẫn linh hoạt trong việc xử lý các nhu cầu mở rộng như báo cáo và tích hợp.

Ngoài ra, việc tách mô hình đọc khỏi dữ liệu giao dịch cho phép tối ưu hóa truy vấn mà không ảnh hưởng đến hiệu năng của hệ thống chính.

\paragraph{Các thành phần chính.}
\begin{itemize}
  \item Service nghiệp vụ
  \item Cơ sở dữ liệu giao dịch
  \item Outbox / hàng đợi sự kiện
  \item Bộ xử lý nền
  \item Mô hình đọc
  \item RabbitMQ và các consumer sự kiện
\end{itemize}

\paragraph{Lý do thiết kế.}
\begin{itemize}
  \item Giữ ranh giới service rõ ràng
  \item Tách xử lý bất đồng bộ để giảm phụ thuộc
  \item Tối ưu hóa truy vấn thông qua mô hình đọc
\end{itemize}

Ưu điểm của kiến trúc này là vừa đảm bảo tính rõ ràng của giao dịch, vừa hỗ trợ mở rộng linh hoạt. Đổi lại, hệ thống cần được thiết kế cẩn thận về hợp đồng API, sự kiện và quyền sở hữu dữ liệu để tránh phụ thuộc chéo giữa các service.

\subsubsection{Sơ đồ triển khai cho runtime Docker local và cô lập lỗi}
\begin{figure}[H]
\centering
\IfFileExists{32_deployment_diagram_high_availability_and_failure_isolation.pdf}{
  \includegraphics[width=0.97\linewidth,height=0.84\textheight,keepaspectratio]{32_deployment_diagram_high_availability_and_failure_isolation.pdf}
}{
  \IfFileExists{32_deployment_diagram_high_availability_and_failure_isolation.png}{
    \includegraphics[width=0.97\linewidth,height=0.84\textheight,keepaspectratio]{32_deployment_diagram_high_availability_and_failure_isolation.png}
  }{
    \fbox{\parbox{0.82\linewidth}{Thiếu file \textsf{32\_deployment\_diagram\_high\_availability\_and\_failure\_isolation.(pdf/png)}.}}
  }
}
\caption{Sơ đồ triển khai cho runtime Docker local và cô lập lỗi}
\label{fig:irms-deployment-ha-supplement}
\end{figure}

Hình~\ref{fig:irms-deployment-ha-supplement} mô tả môi trường runtime thực tế qua \textsf{Docker Compose}, tập trung vào cấu trúc container, luồng sự kiện và cơ chế cô lập lỗi (fault isolation). 

\begin{itemize}
    \item \textbf{Luồng Request:} Request từ \textsf{frontend container} (React/Node.js) đi qua \textsf{api-gateway} (sử dụng \textsf{GatewayProxyController}) để định tuyến REST đến các service Spring Boot. Phạm vi triển khai hiện tại tập trung vào cụm container đơn, chưa bao gồm Load Balancing hay Failover.
    
    \item \textbf{Quản lý dữ liệu:} Dữ liệu lưu tại \textsf{PostgreSQL container} với volume \textsf{postgres-data}. \textsf{irms-migration} sử dụng Flyway để thực hiện migration trước khi khởi động các service nghiệp vụ. Hiện tại hệ thống dùng một instance vật lý duy nhất, chưa triển khai Replication hay Hot-standby.
    
    \item \textbf{Cô lập lỗi (Fault Isolation):} Tách biệt giao dịch lõi và các tác vụ phụ (side-effects). Giao dịch chính được xử lý đồng bộ; các tác vụ như cập nhật Read Model, gửi thông báo và ghi audit log được đẩy qua Outbox Pattern (\textsf{outbox\_events}, \textsf{processed\_events}) và \textsf{RabbitMQ}. Cơ chế này đảm bảo sự cố tại các module hỗ trợ không làm gián đoạn luồng nghiệp vụ chính tại nhà hàng.
    
    \item \textbf{Xử lý kịch bản lỗi:}
    \begin{itemize}
        \item \textsf{Gateway/Service} lỗi: Endpoint tương ứng sẽ ngừng hoạt động dựa trên healthcheck của Docker.
        \item Bộ xử lý nền lỗi: Sự kiện được lưu giữ trong Outbox table hoặc hàng đợi \textsf{RabbitMQ} để retry theo cấu hình.
        \item Side-effect chậm: Không block luồng giao dịch chính; ưu tiên hoàn tất thanh toán và chốt đơn trước.
    \end{itemize}
    
    \item \textbf{Suy giảm dịch vụ (Graceful Degradation):} Kiến trúc cho phép hệ thống duy trì hoạt động phần lõi khi các thành phần phụ gặp sự cố. Việc chấp nhận suy giảm có kiểm soát (chậm báo cáo/thông báo) giúp hệ thống tránh khỏi lỗi dây chuyền (cascading failure).
\end{itemize}
\paragraph{Mô tả}
\begin{itemize}
  \item \textbf{frontend container} tách giao diện React và proxy \textsf{/api} khỏi các service backend.
  \item \textbf{api-gateway} định tuyến REST đến các container service Spring Boot theo cấu hình thật.
  \item \textbf{PostgreSQL container} lưu dữ liệu giao dịch và các nhóm bảng hỗ trợ báo cáo trong một cơ sở dữ liệu vật lý duy nhất.
  \item \textbf{RabbitMQ container} cùng \textsf{outbox\_events} và \textsf{processed\_events} ngăn các tác vụ nền làm nghẽn trực tiếp tuyến giao dịch.
  \item \textbf{Reporting tables and views} cho phép tách nhu cầu đọc báo cáo khỏi phần xử lý lệnh nghiệp vụ.
\end{itemize}



Ngoài khả năng cô lập lỗi, sơ đồ cũng thể hiện rõ hướng mở rộng của hệ thống. Việc bổ sung thêm nút ứng dụng hoặc bộ xử lý nền có thể được thực hiện về sau, nhưng trong tài liệu hiện tại cần phân biệt rõ giữa cấu hình đã có trong \textsf{docker-compose.yml} và phương án mở rộng tương lai.

\subsubsection{Install View cho gói triển khai Docker local}

Deployment View ở trên trả lời các container và node runtime được đặt ở đâu. Install View dưới đây bổ sung góc nhìn còn lại: artifact, cấu hình và migration nào được đóng gói vào từng container khi chạy theo \textsf{Docker Compose}. Phần này chỉ mô tả runtime hiện tại, không bổ sung các hạ tầng như Kubernetes, dịch vụ đám mây, cân bằng tải, bản sao đọc dữ liệu hoặc nhà cung cấp ngoài khi chưa thuộc phạm vi triển khai.

\begin{figure}[H]
\centering
\IfFileExists{37_install_view_docker_local_artifacts.pdf}{
  \includegraphics[width=0.97\linewidth,height=0.84\textheight,keepaspectratio]{37_install_view_docker_local_artifacts.pdf}
}{
  \IfFileExists{37_install_view_docker_local_artifacts.png}{
    \includegraphics[width=0.97\linewidth,height=0.84\textheight,keepaspectratio]{37_install_view_docker_local_artifacts.png}
  }{
    \fbox{\parbox{0.82\linewidth}{Thiếu file \textsf{37\_install\_view\_docker\_local\_artifacts.(pdf/png)}.}}
  }
}

\caption{UML Deployment Diagram cho Install View của gói Docker local}
\label{fig:irms-install-view-docker-local}
\end{figure}

Hình~\ref{fig:irms-install-view-docker-local} dùng UML Deployment Diagram để thể hiện các node runtime và artifact cài đặt của gói Docker local. \textsf{irms-backend:local} chứa \textsf{/app/irms.jar}, các file \textsf{application*.properties} và migration SQL; cùng artifact JAR này được triển khai vào \textsf{api-gateway}, các service Spring Boot và \textsf{irms-migration}, nhưng vai trò runtime khác nhau nhờ profile và biến môi trường trong \textsf{docker-compose.yml}. \textsf{irms-frontend:local} chứa bundle Nitro ở \textsf{.output/server/index.mjs}; \textsf{nitro.config.ts} điều khiển proxy \textsf{/api/**} sang \textsf{api-gateway}. Các node trạng thái chỉ gồm \textsf{PostgreSQL} với \textsf{postgres-data} và \textsf{RabbitMQ} với \textsf{rabbitmq-data}, đúng với cấu hình hiện tại.

\begin{table}[H]
\centering
\footnotesize
\renewcommand{\arraystretch}{1.15}
\caption{Install View cho artifact backend và frontend}
\label{tab:irms-install-view}
\begin{tabularx}{\textwidth}{>{\raggedright\arraybackslash}p{0.19\linewidth} >{\raggedright\arraybackslash}p{0.36\linewidth} >{\raggedright\arraybackslash}X}
\toprule
\textbf{Artifact/config item} & \textbf{Nơi sử dụng} & \textbf{Vai trò runtime} \\
\midrule
Backend Spring Boot artifact \textsf{irms.jar} & \texttt{backend/Dockerfile}; \texttt{backend/pom.xml}; \texttt{docker-compose.yml}. & Image \textsf{irms-backend:local} được dùng lại cho \textsf{api-gateway}, các service nghiệp vụ và \textsf{irms-migration}. Vai trò service được chọn bằng \texttt{SPRING\_PROFILES\_ACTIVE}, \texttt{IRMS\_RUNTIME\_SERVICE} và \texttt{SERVER\_PORT}. \\
Migration runtime và SQL migrations & Giai đoạn khởi tạo schema dữ liệu. & Container \textsf{irms-migration} chạy Flyway trước service nghiệp vụ, với điều kiện \texttt{service\_completed\_successfully}. \\
Frontend bundle / Nitro server & \texttt{frontend/Dockerfile}; \texttt{frontend/package.json}; \texttt{frontend/nitro.config.ts}. & Image \textsf{irms-frontend:local} chạy \texttt{node .output/server/index.mjs}, phục vụ UI React/TanStack và proxy \textsf{/api/**} tới \textsf{api-gateway}. \\
Runtime application properties & \texttt{application.properties}; \texttt{application-docker.properties}; \texttt{application-\char`*-service.properties}. & Cấu hình datasource PostgreSQL, RabbitMQ, Flyway, CORS, security token, timeout phiên, outbox và URL service nội bộ. \\
\bottomrule
\end{tabularx}
\end{table}

\begin{table}[H]
\centering
\scriptsize
\renewcommand{\arraystretch}{1.15}
\setlength{\tabcolsep}{2pt}
\caption{Install View cho thành phần trạng thái và cấu hình Docker}
\label{tab:irms-install-stateful-view}
\begin{tabularx}{\textwidth}{>{\raggedright\arraybackslash}p{0.19\linewidth} >{\raggedright\arraybackslash}p{0.36\linewidth} >{\raggedright\arraybackslash}X}
\toprule
\textbf{Thành phần} & \textbf{Nơi sử dụng} & \textbf{Vai trò runtime} \\
\midrule
\textsf{irms-postgres} và \textsf{postgres-data} & \texttt{docker-compose.yml}; \texttt{docker/postgres/initdb}. & Một PostgreSQL container vật lý duy nhất lưu dữ liệu giao dịch, \textsf{outbox\_events}, \textsf{processed\_events}, các bảng projection báo cáo và dữ liệu identity/audit trong cùng PostgreSQL; các nhóm này không phải database riêng. \\
\textsf{irms-rabbitmq} và \textsf{rabbitmq-data} & \texttt{docker-compose.yml}; \texttt{application.properties}; \texttt{application-docker.properties}. & Broker cho tuyến event-driven; các service dùng \texttt{spring.rabbitmq.*}, manual ack và publisher confirm khi \texttt{IRMS\_RABBITMQ\_ENABLED=true}. \\
Thứ tự khởi động Docker & \texttt{depends\_on} và healthcheck trong \texttt{docker-compose.yml}. & \textsf{postgres} và \textsf{rabbitmq} healthy trước, \textsf{irms-migration} hoàn tất, sau đó Spring Boot services chạy và cuối cùng frontend phụ thuộc \textsf{api-gateway}. \\
\bottomrule
\end{tabularx}
\end{table}

Bảng thể hiện thứ tự cài đặt/khởi động của phạm vi triển khai hiện tại: \textsf{postgres} và \textsf{rabbitmq} phải healthy trước, \textsf{irms-migration} chạy thành công để tạo/cập nhật schema, sau đó các service Spring Boot và cuối cùng là \textsf{frontend} phụ thuộc vào \textsf{api-gateway}. Đây là Install View của runtime Docker local hiện có, không phải mô tả triển khai cloud hay cụm sẵn sàng cao.

\subsection{Nguyên tắc thiết kế}\label{sec:architecture-principles}
Kiến trúc của IRMS bám theo các nguyên tắc thiết kế sau:
\begin{itemize}
  \item \textbf{High Cohesion, Low Coupling:} Mỗi service chỉ expose API, event và data contract cần thiết; ẩn hoàn toàn chi tiết triển khai (encapsulation).
  \item \textbf{Separation of Concerns (SoC):} Tách biệt rõ rệt vai trò giữa Controller, Application Service, Domain và Infrastructure.
  \item \textbf{Dependency Inversion:} Domain core chỉ phụ thuộc vào abstraction (interface); các adapter hạ tầng phụ thuộc ngược lại vào core, tránh gắn chặt vào library hoặc provider cụ thể.
  \item \textbf{Synchronous API (Core Transactions):} Xử lý đồng bộ các nghiệp vụ cần phản hồi tức thời như đặt bàn, gọi món và thanh toán để đảm bảo tính nhất quán trạng thái.
  \item \textbf{Asynchronous Events (Side-effects):} Các tác vụ phụ như cảnh báo kho, cập nhật báo cáo, thông báo và audit log được xử lý bất đồng bộ qua Outbox Pattern và Event Bus.
  \item \textbf{Domain-centric Modeling:} Tập trung mô hình hóa hành vi nghiệp vụ thực tế (Domain Behavior) thay vì thiết kế dựa trên quan hệ dữ liệu (Database-centric) thuần túy.
\end{itemize}

\subsection{Áp dụng nguyên lý SOLID vào thiết kế}\label{sec:solid-design}

Phần này đối chiếu trực tiếp các nguyên tắc SOLID với thiết kế thực tế (implementation) của IRMS qua các class và interface tiêu biểu, thay vì liệt kê lý thuyết đơn thuần. Các interface như \textsf{PaymentGateway}, \textsf{NotificationChannelAdapter} và \textsf{ReceiptDeliveryAdapter} cụ thể hóa quyết định dùng boundary adapters trong \adrref{ADR-0011}{adr:0011-boundary-adapters}; cách tách lớp \texttt{api}, \texttt{application}, \texttt{domain} và \texttt{infrastructure} liên kết với \adrref{ADR-0013}{adr:0013-service-layering}.

\begin{table}[H]
\centering
\footnotesize
\renewcommand{\arraystretch}{1.15}
\caption{Đối chiếu nguyên tắc SOLID với thiết kế IRMS}
\label{tab:solid-thiết kế-evidence}
\begin{tabularx}{\textwidth}{>{\raggedright\arraybackslash}p{0.13\linewidth} >{\raggedright\arraybackslash}X >{\raggedright\arraybackslash}X}
\toprule
\textbf{Nguyên tắc} & \textbf{Ví dụ thiết kế tiêu biểu} & \textbf{Ý nghĩa thiết kế} \\
\midrule
SRP & \texttt{BillingPaymentService}; \texttt{BillingReceiptService}; \texttt{KitchenTicketTransitionPolicy}; \texttt{ManualAckConsumerSupport}; \texttt{RabbitMqOutboxRelay}. & Trách nhiệm được tách theo service, policy, adapter và consumer; khi thay đổi thanh toán, trạng thái bếp hoặc retry RabbitMQ, phạm vi sửa được khoanh vào nhóm class đúng vai trò. \\
OCP & \texttt{PaymentGateway} với \texttt{CashPaymentGateway}, \texttt{ExternalReferencePaymentGateway}; \texttt{BillSplitStrategy} với các strategy chia hóa đơn; \texttt{ReorderRecommendationStrategy}. & Các điểm dễ thay đổi được mở qua interface/strategy/policy, nên có thể thêm biến thể mới bằng class hiện thực mới và cấu hình bean. \\
LSP & \texttt{BillingPaymentService} nhận \texttt{List<PaymentGateway>} và gọi \texttt{supports(...)}; \texttt{BillingSplitService} nhận \texttt{List<BillSplitStrategy>}. & Service sử dụng contract ổn định; implementation có thể thay thế trong phạm vi hợp đồng đã định nghĩa. \\
ISP & \texttt{PaymentGateway}; \texttt{ReceiptDeliveryAdapter}; \texttt{NotificationChannelAdapter}; \texttt{KitchenTicketItemRepositoryPort}; \texttt{DomainEventPublisher}. & Interface được giữ nhỏ và gắn với một vai trò cụ thể, giúp client chỉ phụ thuộc vào hành vi nó thật sự cần. \\
DIP & \texttt{BillingPaymentService} phụ thuộc \texttt{BillPaymentRepository}, \texttt{PaymentGateway}, \texttt{DomainEventPublisher}; \texttt{BillingReceiptService} phụ thuộc \texttt{ReceiptDeliveryAdapter}; \texttt{JdbcKitchenTicketItemRepository} implements \texttt{KitchenTicketItemRepositoryPort}. & Service cấp cao phụ thuộc abstraction/port, còn hạ tầng hiện thực chi tiết JDBC, RabbitMQ hoặc adapter kênh. \\
\bottomrule
\end{tabularx}
\end{table}

\paragraph{Thanh toán và biên nhận.}
\textsf{BillingPaymentService} không biết chi tiết xử lý từng phương thức thanh toán, mà làm việc qua \textsf{PaymentGateway}. Các hiện thực như \textsf{CashPaymentGateway} và \textsf{ExternalReferencePaymentGateway} được xem là adapter thanh toán ở biên. Tương tự, \textsf{BillingReceiptService} phát hành biên nhận qua \textsf{ReceiptDeliveryAdapter} với \textsf{DigitalReceiptDeliveryAdapter} và \textsf{PrintReceiptDeliveryAdapter}. Đây là ví dụ mạnh cho OCP, ISP và DIP: service lõi giữ luồng thanh toán, còn biến thể phương thức/kênh nằm ở lớp hiện thực.

\paragraph{Thông báo.}
Nhóm \textsf{NotificationChannelAdapter} được tách thành các kênh nhỏ, gồm:
\begin{itemize}
  \item \textsf{EmailNotificationChannelAdapter}, \textsf{SmsNotificationChannelAdapter};
  \item \textsf{InAppNotificationChannelAdapter}, \textsf{PrintNotificationChannelAdapter};
  \item \textsf{PhoneNotificationChannelAdapter}.
\end{itemize}
Những class này giúp service thông báo kiểm tra và ghi nhận kênh gửi mà không biến phần notification thành một interface quá lớn. Các adapter này được mô tả như điểm mở rộng nội bộ cho các kênh gửi, không làm lõi thông báo phụ thuộc trực tiếp vào một kênh cụ thể.

\paragraph{Policy và strategy nghiệp vụ.}
\textsf{DomainPolicyConfiguration} đăng ký các nhóm policy/strategy như:
\begin{itemize}
  \item \textsf{OrderStatusPolicy}, \textsf{ComboPricingPolicy};
  \item \textsf{KitchenTicketTransitionPolicy}, \textsf{KitchenDeadlinePolicy};
  \item \textsf{LowStockSeverityPolicy}, \textsf{ReorderRecommendationStrategy};
  \item các hiện thực \textsf{BillSplitStrategy}.
\end{itemize}
Cách tách này làm rõ SRP: quy tắc chuyển trạng thái, tính deadline, đề xuất nhập hàng hoặc tách hóa đơn không bị trộn vào controller. Đồng thời, các strategy như \textsf{BillSplitStrategy} và \textsf{ReorderRecommendationStrategy} là điểm mở rộng OCP được thể hiện rõ trong thiết kế.

\paragraph{Outbox, RabbitMQ và bộ xử lý sự kiện nền.}
Ở tuyến event-driven, \textsf{TransactionalOutboxPublisher} hiện thực \textsf{DomainEventPublisher}, \textsf{RabbitMqOutboxRelay} phụ trách chuyển outbox sang RabbitMQ, còn \textsf{ManualAckConsumerSupport} gom logic ack, retry, deduplication và dead-letter cho consumer. Phần này là ví dụ SRP và DIP ở tầng hạ tầng: service nghiệp vụ chỉ publish event qua abstraction, còn cách phát, retry hoặc dead-letter nằm trong component kỹ thuật riêng.

\paragraph{Giới hạn cần tránh diễn giải quá mức.}
Các ví dụ trên cho thấy IRMS có áp dụng SOLID tại những điểm thiết kế chính, đặc biệt ở service, adapter, policy, strategy, port truy cập dữ liệu và event infrastructure. Điều đó không có nghĩa mọi class đều hoàn hảo theo SOLID. Ngoài ra, các interface như \textsf{PaymentGateway}, \textsf{ReceiptDeliveryAdapter} và \textsf{NotificationChannelAdapter} xác định ranh giới trừu tượng, không làm cho lõi nghiệp vụ phụ thuộc vào một kênh thanh toán, thông báo hoặc in ấn cụ thể.

\subsection{Khả năng mở rộng trong tương lai}
\begin{figure}[H]
\centering
\IfFileExists{36_future_extensibility_evolution_path.pdf}{
  \includegraphics[width=\textwidth,height=0.82\textheight,keepaspectratio]{36_future_extensibility_evolution_path.pdf}
}{
  \IfFileExists{36_future_extensibility_evolution_path.png}{
    \includegraphics[width=\textwidth,height=0.82\textheight,keepaspectratio]{36_future_extensibility_evolution_path.png}
  }{
    \fbox{\parbox{0.82\linewidth}{Thiếu file \textsf{36\_future\_extensibility\_evolution\_path.(pdf/png)}.}}
  }
}
\caption{Lộ trình mở rộng kiến trúc của IRMS trong tương lai}
\label{fig:irms-future-extensibility}
\end{figure}

Hình~\ref{fig:irms-future-extensibility} mô tả lộ trình tiến hóa (evolutionary path) của hệ thống khi mở rộng quy mô (tăng thiết bị, chi nhánh, tích hợp) mà không cần tái cấu trúc tầng lõi (Core Domain). 

\paragraph{Trạng thái kiến trúc hiện tại:}
IRMS áp dụng mô hình SOA kết hợp Event-driven. Các nghiệp vụ chính (Core Transactions) sử dụng Sync API để đảm bảo tính nhất quán; các tác vụ phụ (Side-effects) được xử lý bất đồng bộ qua Event. Thiết kế này cho phép các service hỗ trợ như \textsf{Reporting}, \textsf{Notification}, \textsf{Inventory} có thể scale độc lập tùy theo tải trọng thực tế.

\paragraph{Nguyên tắc mở rộng không phá vỡ kiến trúc:}
Khả năng mở rộng (Extensibility) của IRMS được đo lường bằng việc duy trì sự ổn định của Core Domain (Order, Billing) khi thay đổi các thành phần biên:
\begin{itemize}
    \item \textbf{Tích hợp thanh toán/thông báo:} Chỉ tác động đến \textsf{PaymentGateway} và \textsf{NotificationChannelAdapter}.
    \item \textbf{Mở rộng quy mô:} Thêm chi nhánh hoặc thiết bị (Kiosk, Mobile) chỉ yêu cầu điều chỉnh tại tầng Deployment và API layer.
    \item \textbf{Điểm mở rộng chiến lược:} Các thành phần như Gateway, Adapter, Policy (quy tắc định giá/tồn kho) và Read Model cần được giữ ổn định để đảm bảo hệ thống tiến hóa có kiểm soát.
\end{itemize}

\paragraph{Các hướng mở rộng chính:}
\begin{itemize}
    \item \textbf{Mở rộng tích hợp (Integration):} Thêm Provider (thanh toán, thông báo) hoặc thiết bị ngoại vi qua mô hình Ports \& Adapters.
    \item \textbf{Mở rộng theo tải (Scaling):} Tăng số lượng instance ứng dụng, Background Worker hoặc tách biệt hoàn toàn Read Model để tối ưu truy vấn.
    \item \textbf{Mở rộng phạm vi vận hành:} Hỗ trợ đa chi nhánh và quản trị tập trung (Centralized Management) trên cùng một mô hình nghiệp vụ lõi.
    \item \textbf{Mở rộng theo dịch vụ:} Chỉ tách thành các Microservices độc lập đối với những vùng có nhịp thay đổi (rate of change) hoặc tải trọng khác biệt hoàn toàn.
\end{itemize}

\section{Thiết kế chi tiết}
